% Luc Destombe --- licence CC BY-NC-SA 4.0
% https://creativecommons.org/licenses/by-nc-sa/4.0/deed.fr
% Fichier autonome : compiler avec pdflatex, deux fois (signets, nombre de pages).
\documentclass[11pt,a4paper]{article}
\makeatletter
\newif\iflycee@unecolonne
\lycee@unecolonnetrue
\newif\iflycee@palatino
\lycee@palatinotrue

\RequirePackage[utf8]{inputenc}
\RequirePackage[T1]{fontenc}
\RequirePackage[french]{babel}
\RequirePackage{etoolbox}

\newcommand{\ShorthandsOff}{\@ifundefined{shorthandoff}{}{\shorthandoff{;:!?}}}
\newcommand{\ShorthandsOn}{\@ifundefined{shorthandon}{}{\shorthandon{;:!?}}}
\ShorthandsOff
\AtBeginDocument{\ShorthandsOn}
\AtBeginEnvironment{tikzpicture}{\ShorthandsOff}

\RequirePackage{geometry}
\RequirePackage{fancyhdr}
\RequirePackage{lastpage}
\RequirePackage{multicol}
\RequirePackage{array}
\RequirePackage{enumitem}
\RequirePackage{xcolor}
\RequirePackage{needspace}

\RequirePackage{amsmath,amssymb}
\RequirePackage{mathpazo}
\DeclareMathSizes{10.95}{10.95}{8}{5.5}
\begingroup\catcode`\_=\active \catcode`\^=\active
\gdef_#1{\sb{\scriptscriptstyle #1}}
\gdef^#1{\sp{\scriptscriptstyle\lycee@moinsepais #1}}
\endgroup
\newcommand\lycee@moins{\mathbin{%
  \setbox\z@\hbox{$\m@th\scriptscriptstyle\lycee@moinsorig$}%
  \dimen@=.55\wd\z@ \advance\dimen@-.5pt
  \hbox to.75\wd\z@{\hss$\m@th\scriptscriptstyle\vcenter{\hbox{%
    \tikz\draw[line width=.5pt,line cap=round](0,0)--(\the\dimen@,0);}}$\hss}}}
\begingroup\lccode`\~=`\- \lowercase{\endgroup\let~}\lycee@moins
\newcommand\lycee@moinsepais{\mathcode`\-="8000 }
\AtBeginDocument{\mathchardef\lycee@moinsorig=\mathcode`\-\relax}
\AtBeginDocument{\catcode`\_=12 \mathcode`\_="8000
  \catcode`\^=12 \mathcode`\^="8000 }
\newcommand\lycee@exposants{%
  \fontdimen13\textfont2=.48\fontdimen6\textfont2
  \fontdimen14\textfont2=.48\fontdimen6\textfont2
  \fontdimen15\textfont2=.48\fontdimen6\textfont2
  \fontdimen13\scriptfont2=.48\fontdimen6\scriptfont2
  \fontdimen14\scriptfont2=.48\fontdimen6\scriptfont2
  \fontdimen15\scriptfont2=.48\fontdimen6\scriptfont2}
\every@math@size\expandafter{\the\every@math@size\lycee@exposants}
\newlength{\retraitcalculs}
\setlength{\retraitcalculs}{2em}

\RequirePackage{tikz}
\usetikzlibrary{arrows.meta,calc,babel,shapes.misc}
\RequirePackage[most]{tcolorbox}
\tcbuselibrary{skins,breakable}

\definecolor{coulDef}{HTML}{1F4E79}
\definecolor{coulProp}{HTML}{7030A0}
\definecolor{coulRem}{HTML}{C55A11}
\definecolor{coulEx}{HTML}{2E7D32}
\definecolor{coulExo}{HTML}{B03A2E}
\definecolor{coulPart}{HTML}{1F4E79}
\definecolor{coulMeth}{HTML}{1B6E4F}
\definecolor{coulCourbe}{HTML}{0B5ED7}
\definecolor{coulCourbeB}{HTML}{C00000}
\definecolor{coulCourbeC}{HTML}{2E7D32}
\definecolor{coulGrille}{HTML}{8C8C8C}
\definecolor{coulRep}{HTML}{1B6E4F}
\definecolor{coulBar}{HTML}{2E7D32}

\newcommand{\R}{\mathbb{R}}
\newcommand{\Rs}{\mathbb{R}^{*}}
\renewcommand{\le}{\leqslant}
\renewcommand{\ge}{\geqslant}
\renewcommand{\approx}{\simeq}
\newcommand{\vect}[1]{\overrightarrow{#1}}
\newcommand{\norme}[1]{\left\lVert #1 \right\rVert}
\newcommand{\intervalle}[2]{\left[#1\,;#2\right]}
\RequirePackage{cancel}
\renewcommand{\CancelColor}{\color{coulExo}}
\newcommand{\fois}[1]{\times{\color{coulExo}#1}}
\newcommand{\barre}[1]{\cancel{#1}}

\tikzset{grille/.style={coulGrille,line width=0.4pt}}
\newcommand{\quadrillage}[4]{\draw[grille,step=1] (#1,#2) grid (#3,#4);}
\tikzset{
  croix/.style={cross out,draw,line width=0.6pt,inner sep=0pt,minimum size=4pt},
  croixdroite/.style={croix,rotate=45,line width=0.8pt},
}
\newcommand{\pt}[3]{\path (#1) node[croix]{} node[#2,font=\small,inner sep=2.5pt]{$#3$};}

\newcommand{\lycee@repere}[4]{%
  \draw[grille,step=1] (#1,#2) grid (#3,#4);
  \draw[-{Stealth[length=2mm]},thick] (#1,0)--({#3+0.4},0) node[below left=-1pt]{$x$};
  \draw[-{Stealth[length=2mm]},thick] (0,#2)--(0,{#4+0.4}) node[below left=-1pt]{$y$};
  \node[below left,font=\scriptsize,inner sep=1.5pt] at (0,0) {$0$};}
\newcommand{\lycee@gradx}[1]{\foreach \k/\l in {#1}{%
  \draw (\k,0.07)--(\k,-0.07) node[below,font=\scriptsize,inner sep=1.5pt]{$\l$};}}
\newcommand{\lycee@grady}[1]{\foreach \k/\l in {#1}{%
  \draw (0.07,\k)--(-0.07,\k) node[left,font=\scriptsize,inner sep=1.5pt]{$\l$};}}
\newcommand{\lycee@intff}[2]{\left[#1\,;#2\right]}
\newcommand{\lycee@intfo}[2]{\left[#1\,;#2\right[}
\newcommand{\lycee@intof}[2]{\left]#1\,;#2\right]}
\newcommand{\lycee@intoo}[2]{\left]#1\,;#2\right[}
\newcommand{\lycee@ens}[1]{\left\{#1\right\}}
\AtBeginDocument{%
  \providecommand{\repere}{\lycee@repere}%
  \providecommand{\gradx}{\lycee@gradx}%
  \providecommand{\grady}{\lycee@grady}%
  \providecommand{\Cf}{\mathcal{C}\sb{\scriptscriptstyle f}}%
  \providecommand{\Cg}{\mathcal{C}\sb{\scriptscriptstyle g}}%
  \providecommand{\intff}{\lycee@intff}%
  \providecommand{\intfo}{\lycee@intfo}%
  \providecommand{\intof}{\lycee@intof}%
  \providecommand{\intoo}{\lycee@intoo}%
  \providecommand{\ens}{\lycee@ens}}

\newlist{questions}{enumerate}{3}
\setlist[questions,1]{label=\arabic*\textdegree),leftmargin=*,itemsep=2pt,topsep=3pt}
\setlist[questions,2]{label=\alph*),leftmargin=*,itemsep=1pt,topsep=2pt}
\setlist[questions,3]{label=\roman*),leftmargin=*,itemsep=1pt,topsep=1pt}
\setlist[itemize]{label=\textbullet,leftmargin=*,itemsep=2pt,topsep=3pt}

\newcommand{\besoinplace}[1]{\par\penalty\@M\begingroup
  \dimen@=#1\relax\dimen@ii=\pagegoal\advance\dimen@ii-\pagetotal
  \ifdim\dimen@>\dimen@ii\ifdim\dimen@ii>\z@\vfil\fi\break\fi\endgroup}

\newcommand{\lycee@pied}{\thepage/\pageref{LastPage}}
\newcommand{\entete}[2]{%
  \pagestyle{fancy}\fancyhf{}%
  \fancyhead[L]{\footnotesize\color{coulPart}\textbf{#1}}%
  \fancyhead[R]{\footnotesize\color{coulPart}\textbf{#2}}%
  \fancyfoot[C]{\footnotesize\lycee@pied}%
  \renewcommand{\headrulewidth}{0.4pt}%
  \renewcommand{\headrule}{\hbox to\headwidth{%
    \color{coulPart!40}\leaders\hrule height \headrulewidth\hfill}}}

\RequirePackage{ccicons}
\newcommand{\lycee@licence}{{\scriptsize\color{gray}%
  \copyright\ Luc Destombe --- \ccbyncsa\ CC BY-NC-SA 4.0}}
\newif\iflycee@licence \lycee@licencetrue
\newcommand{\sanslicence}{\lycee@licencefalse}
\AtBeginDocument{\iflycee@licence\AddToHookNext{shipout/foreground}{%
  \put(0,0){\rlap{\kern\dimexpr1in+\hoffset+\oddsidemargin+\textwidth\relax
    \llap{\raisebox{-\dimexpr1in+\voffset+\topmargin+\headheight+\headsep
      +\textheight+\footskip\relax}{\lycee@licence}}}}}\fi}

\newcommand{\formulesserrees}{%
  \setlength{\abovedisplayskip}{3pt}\setlength{\belowdisplayskip}{3pt}%
  \setlength{\abovedisplayshortskip}{2pt}\setlength{\belowdisplayshortskip}{3pt}}

\setlength{\parindent}{0pt}

\RequirePackage[implicit=false,hidelinks,pdfencoding=auto,psdextra,bookmarksopen,
  bookmarksopenlevel=1,pdfcreator={},pdfproducer={}]{hyperref}
\RequirePackage{bookmark}
\hypersetup{pdfauthor={Luc Destombe},pdfkeywords={CC BY-NC-SA 4.0}}
\newcounter{lycee@signet}
\newcommand{\signet}[3][]{\stepcounter{lycee@signet}%
  \edef\lycee@dest{\if\relax\detokenize{#1}\relax
    signet.\arabic{lycee@signet}\else#1\fi}%
  \hypertarget{\lycee@dest}{}\bookmark[level=#2,dest=\lycee@dest]{#3}}
\pdfstringdefDisableCommands{%
  \def\R{R}\def\vect#1{#1}\def\Cf{Cf}\def\Cg{Cg}\def\fois#1{×#1}%
  \def\barre#1{#1}\def\intervalle#1#2{[#1;#2]}\def\intff#1#2{[#1;#2]}%
  \def\textsuperscript#1{#1}\def\dfrac#1#2{#1/#2}\def\frac#1#2{#1/#2}%
  \def\vec#1{#1⃗}}%
\RequirePackage[official]{eurosym}

\tcbset{exercice corrige/.style={
  enhanced,breakable,before upper=\raggedright,
  colback=white,colframe=coulExo,
  boxrule=0.8pt,arc=2pt,
  left=6pt,right=6pt,top=8pt,bottom=5pt,
  before skip=9pt,after skip=4pt,
  coltitle=white,fonttitle=\bfseries\small,
  attach boxed title to top left={xshift=8pt,yshift=-\tcboxedtitleheight/2},
  boxed title style={colback=coulExo,arc=2pt,boxrule=0pt}}}
\newcounter{exo}
\newtcolorbox{exercice}[1][]{exercice corrige,
  phantom={\signet[exercice-\arabic{exo}]{2}{Exercice \arabic{exo}}},
  title={Exercice \theexo\ifx\relax#1\relax\else{} --- #1\fi}}
\newenvironment{exo}[1][\relax]{\refstepcounter{exo}\begin{exercice}[#1]}{\end{exercice}}

\geometry{top=1.8cm,bottom=1cm,left=1cm,right=1cm,headheight=14pt,footskip=0.6cm}
\renewcommand{\lycee@pied}{\thepage}

\renewcommand{\theexo}{\ifnum\value{exo}<10 0\fi\arabic{exo}}
\tcbset{exercice corrige/.style={
  enhanced,breakable,before upper=\raggedright,
  colback=white,colframe=coulExo,
  boxrule=0.9pt,arc=2pt,
  left=8pt,right=8pt,top=8pt,bottom=6pt,
  before skip=8pt,after skip=6pt,
  coltitle=white,fonttitle=\bfseries,
  attach boxed title to top left={xshift=8pt,yshift=-\tcboxedtitleheight/2},
  boxed title style={colback=coulExo,arc=2pt,boxrule=0pt}}}

\newcommand{\entetefiche}[2]{%
  \begin{center}
    \begin{tcolorbox}[enhanced,width=0.92\linewidth,colback=white,colframe=coulPart,
        boxrule=1.4pt,arc=3pt,halign=center,top=6pt,bottom=6pt]
      {\large\bfseries\color{coulPart} CORRIGÉ --- FICHE D'EXERCICES}\\[3pt]
      {\normalsize\bfseries #1}\\[2pt]
      {\normalsize #2}
    \end{tcolorbox}
  \end{center}}

\newcounter{partie}
\newcommand{\partie}[1]{%
  \stepcounter{partie}%
  \besoinplace{8\baselineskip}\vspace{4pt}%
  \begin{tcolorbox}[enhanced,colback=coulPart!8,colframe=coulPart,boxrule=0pt,
      leftrule=3pt,arc=0pt,left=6pt,right=4pt,top=3pt,bottom=3pt,
      before skip=6pt,after skip=2pt]
    \bfseries\color{coulPart}\Roman{partie} --- #1\signet{1}{\Roman{partie} --- #1}
  \end{tcolorbox}}
\newcommand{\partiesansnum}[1]{%
  \besoinplace{8\baselineskip}\vspace{4pt}%
  \begin{tcolorbox}[enhanced,colback=coulPart!8,colframe=coulPart,boxrule=0pt,
      leftrule=3pt,arc=0pt,left=6pt,right=4pt,top=3pt,bottom=3pt,
      before skip=6pt,after skip=2pt]
    \bfseries\color{coulPart}#1\signet{1}{#1}
  \end{tcolorbox}}

\setlist[questions,1]{label=\arabic*\textdegree),leftmargin=*,itemsep=3pt,topsep=3pt}
\setlist[questions,2]{label=\alph*),leftmargin=*,itemsep=2pt,topsep=2pt}
\setlist[questions,3]{label=\roman*),leftmargin=*,itemsep=1pt,topsep=1pt}
\newlist{expr}{enumerate}{1}
\setlist[expr]{label={\Alph*\ $=$},leftmargin=*,itemsep=2pt,topsep=3pt}

\newcommand{\res}[1]{{\color{coulRep}#1}}
\newcommand{\rem}[1]{\par\smallskip{\small\itshape\color{black!60}$\rightarrow$ #1}}
\newenvironment{calculs}{\par\smallskip\noindent$\displaystyle\begin{aligned}[t]}%
  {\end{aligned}$\par\smallskip}
\newcommand{\justif}[1]{\quad\text{\small\itshape\color{black!60}#1}}
\newcommand{\cas}[1]{\par\smallskip\noindent\textbf{\color{coulExo}#1}\nobreak}
\newcommand{\explique}[2]{\par\smallskip\noindent
  \begin{minipage}[t]{0.6\linewidth}\raggedright #1\end{minipage}\hfill
  \begin{minipage}[t]{0.37\linewidth}\raggedright\leavevmode
    {\small\itshape\color{black!60}#2}\end{minipage}\par}

\setlength{\parskip}{2pt}

\makeatother
\geometry{top=1.8cm,bottom=1.6cm,left=1.8cm,right=1.8cm,headheight=14pt,footskip=30pt}
\entete{Chapitre 7 --- Suites arithmétiques et géométriques --- corrigé}{Première spécialité}
\setlength{\parindent}{0pt}
\setlength{\jot}{4pt}

\let\casligne\cas
\renewcommand{\cas}[1]{\casligne{#1}\par\nopagebreak}

\newtcblisting{python}{%
  listing only,enhanced,
  colback=black!3,colframe=black!35,boxrule=0.5pt,arc=1pt,
  left=4pt,right=4pt,top=1pt,bottom=1pt,
  listing options={language=Python,basicstyle=\ttfamily\small,
    keywordstyle=\color{coulMeth}\bfseries,columns=fullflexible,
    keepspaces=true,showstringspaces=false,aboveskip=0pt,belowskip=0pt}}
\newcommand{\reperesuite}[4]{%
  \draw[grille,step=1] (-1,#2) grid (#1,#3);
  \draw[-{Stealth[length=2mm]},thick] (-1,0)--({#1+0.4},0)
    node[below left=-1pt,font=\footnotesize]{$n$};
  \draw[-{Stealth[length=2mm]},thick] (0,#2)--(0,{#3+0.4})
    node[below left=-1pt,font=\footnotesize]{$#4$};
  \node[below left,font=\scriptsize,inner sep=1.5pt] at (0,0) {$0$};}

\begin{document}

\entetefiche{Chapitre 7 --- Suites arithmétiques et géométriques}{Première spécialité mathématiques}

\partie{Suites arithmétiques}

\begin{exo}[Reconnaître une suite arithmétique]
\explique{Pour chaque début de suite, on calcule l'écart entre deux termes qui se
suivent : le terme suivant moins le précédent.}{une suite arithmétique ajoute toujours
le même nombre $r$}
\cas{a)}
\explique{$11-7=4$ ;\; $15-11=4$ ;\; $19-15=4$.}{écart entre deux termes qui se suivent}
\explique{L'écart est toujours $4$ : la suite peut être arithmétique, de raison
$\res{r=4}$.}{on ajoute $4$ à chaque fois}
\cas{b)}
\explique{$24-30=-6$ ;\; $18-24=-6$ ;\; $12-18=-6$ ;\; $6-12=-6$.}{toujours dans le même
ordre : suivant moins précédent}
\explique{L'écart est toujours $-6$ : la suite peut être arithmétique, de raison
$\res{r=-6}$.}{on enlève $6$ : la raison est négative}
\rem{Erreur fréquente : calculer $30-24=6$ et donner $r=6$. Une suite qui diminue a une
raison négative.}
\cas{c)}
\explique{$2-1=1$ mais $4-2=2$.}{deux écarts différents suffisent}
\explique{L'écart n'est pas constant : la suite \res{n'est pas arithmétique}.}{inutile de
calculer les autres écarts}
\cas{d)}
\explique{$3-2{,}5=0{,}5$ ;\; $3{,}5-3=0{,}5$ ;\; $4-3{,}5=0{,}5$.}{la raison peut être
un nombre décimal}
\explique{L'écart est toujours $0{,}5$ : la suite peut être arithmétique, de raison
$\res{r=0{,}5}$.}{on ajoute $0{,}5$ à chaque fois}
\rem{On dit \og peut être\fg{} : on ne connaît que les premiers termes, rien ne garantit
que l'écart reste le même ensuite.}
\end{exo}

\begin{exo}[Traduire une situation]
\explique{Dans chaque cas, on cherche la valeur de départ ($u_{0}$), puis ce qui se passe
quand on avance d'un pas (de $n$ à $n+1$).}{la phrase de l'énoncé dit ce qu'on ajoute}
\cas{a)}
\explique{Pour $0$~km, on ne paie que la prise en charge : $u_{0}=4$.}{$u_{0}$ : la
situation de départ}
\explique{Chaque kilomètre de plus coûte $1{,}5$~€ : $u_{n+1}=u_{n}+1{,}5$.}{prix pour
$n+1$~km = prix pour $n$~km $+\,1{,}5$}
\explique{On ajoute toujours $1{,}5$ : la suite est arithmétique, de raison
$\res{r=1{,}5}$.}{définition}
\cas{b)}
\explique{Au départ, la bougie mesure $18$~cm : $u_{0}=18$.}{au bout de $0$ heure}
\explique{Chaque heure, elle perd $1{,}2$~cm : $u_{n+1}=u_{n}-1{,}2$.}{\og raccourcit\fg{} :
on enlève}
\explique{La suite est arithmétique, de raison $\res{r=-1{,}2}$.}{enlever $1{,}2$, c'est
ajouter $-1{,}2$}
\cas{c)}
\explique{On note $u_{n}$ la somme sur le livret au bout de $n$ mois : $u_{0}=300$ et
$u_{n+1}=u_{n}+25$.}{l'énoncé ne nomme pas la suite : on le fait}
\explique{La suite est arithmétique, de raison $\res{r=25}$.}{on ajoute $25$~€ chaque mois}
\end{exo}

\begin{exo}[Calculer les premiers termes]
\explique{Dans une suite arithmétique, on obtient chaque terme en ajoutant la raison au
précédent : $u_{n+1}=u_{n}+r$.}{définition}
\cas{a)}
\begin{calculs}
u_{1} &= 4+3 = \res{7} &&\justif{$u_{1}=u_{0}+r$}\\
u_{2} &= 7+3 = \res{10} &&\justif{$u_{2}=u_{1}+r$}\\
u_{3} &= 10+3 = \res{13} &&\\
u_{4} &= 13+3 = \res{16} &&
\end{calculs}
\cas{b)}
\begin{calculs}
v_{1} &= 10-2{,}5 = \res{7{,}5} &&\justif{ajouter $-2{,}5$, c'est enlever $2{,}5$}\\
v_{2} &= 7{,}5-2{,}5 = \res{5} &&\\
v_{3} &= 5-2{,}5 = \res{2{,}5} &&\\
v_{4} &= 2{,}5-2{,}5 = \res{0} &&
\end{calculs}
\cas{c)}
\explique{Le premier terme est $w_{1}$ : les quatre termes qui suivent sont $w_{2}$,
$w_{3}$, $w_{4}$ et $w_{5}$.}{on ne commence pas toujours au rang $0$}
\begin{calculs}
w_{2} &= -7+6 = \res{-1} &&\justif{$w_{2}=w_{1}+r$}\\
w_{3} &= -1+6 = \res{5} &&\\
w_{4} &= 5+6 = \res{11} &&\\
w_{5} &= 11+6 = \res{17} &&
\end{calculs}
\end{exo}

\begin{exo}[Trois termes à compléter]
\explique{Pour trois termes consécutifs d'une suite arithmétique, l'écart entre le
deuxième et le premier est égal à l'écart entre le troisième et le deuxième.}{les deux
écarts valent $r$}
\begin{calculs}
(2x+1)-(x-1) &= (4x-3)-(2x+1) &&\justif{même écart}\\
2x+1-x+1 &= 4x-3-2x-1 &&\justif{le signe $-$ devant la parenthèse change les signes}\\
x+2 &= 2x-4 &&\justif{on réduit chaque membre}\\
6 &= x &&\justif{on ajoute $4$ et on retire $x$ des deux côtés}
\end{calculs}
\explique{Avec $x=6$ : $x-1=5$ ;\; $2x+1=13$ ;\; $4x-3=21$. Les trois termes sont
$\res{5}$ ; $\res{13}$ ; $\res{21}$.}{on remplace $x$ par $6$}
\explique{Vérification : $13-5=8$ et $21-13=8$. L'écart est constant : la raison est
$\res{r=8}$.}{on contrôle toujours le résultat}
\rem{Erreur fréquente : oublier les parenthèses et écrire $2x+1-x-1$ : le $-1$ devient
$+1$ quand on retire $(x-1)$.}
\end{exo}

\partie{Terme général d'une suite arithmétique}

\begin{exo}[Exprimer $u_{n}$ en fonction de $n$]
\explique{Pour aller de $u_{p}$ à $u_{n}$, on ajoute $(n-p)$ fois la raison :
$u_{n}=u_{p}+(n-p)\,r$. Avec $p=0$, on retrouve $u_{n}=u_{0}+n\,r$.}{propriété 1 : on
part du terme connu}
\cas{a)}
\begin{calculs}
u_{n} &= u_{0}+n\,r &&\justif{la formule : le premier terme est $u_{0}$}\\
u_{n} &= \res{5+4n} &&\justif{$u_{0}=5$ et $r=4$}\\
u_{30} &= 5+4\times 30 &&\justif{on remplace $n$ par $30$}\\
u_{30} &= \res{125} &&
\end{calculs}
\explique{Contrôle : de $u_{0}$ à $u_{30}$, on ajoute $30$ fois $4$, soit $120$, et
$5+120=125$.}{$30$ sauts de $r$}
\cas{b)}
\explique{$v_{n+1}=v_{n}-3$ : on enlève $3$ à chaque rang, la suite est arithmétique de
raison $r=-3$.}{on lit la raison dans la relation}
\begin{calculs}
v_{n} &= v_{1}+(n-1)\,r &&\justif{le premier terme est $v_{1}$ : $p=1$}\\
v_{n} &= 12-3(n-1) &&\justif{$v_{1}=12$ et $r=-3$}\\
v_{n} &= 12-3n+3 &&\justif{on développe : $-3\times(-1)=+3$}\\
v_{n} &= \res{15-3n} &&\\
v_{40} &= 15-3\times 40 &&\justif{on remplace $n$ par $40$}\\
v_{40} &= \res{-105} &&
\end{calculs}
\explique{Contrôle : $v_{1}=15-3=12$.}{on retrouve la donnée}
\rem{Pourquoi $n-1$ : de $v_{1}$ à $v_{n}$, il y a $n-1$ sauts de $r$ (de $v_{1}$ à
$v_{3}$, deux sauts). Erreur fréquente : écrire $v_{n}=12-3n$, comme si le premier
terme était $v_{0}$.}
\cas{c)}
\begin{calculs}
w_{n} &= w_{5}+(n-5)\,r &&\justif{on part de $w_{5}$ : $p=5$}\\
w_{n} &= 20+2(n-5) &&\justif{$w_{5}=20$ et $r=2$}\\
w_{n} &= 20+2n-10 &&\justif{on développe}\\
w_{n} &= \res{2n+10} &&\\
w_{25} &= 2\times 25+10 &&\\
w_{25} &= \res{60} &&
\end{calculs}
\explique{Contrôle : de $w_{5}$ à $w_{25}$, il y a $20$ sauts de $2$, et
$20+40=60$.}{$n-p=25-5$}
\end{exo}

\begin{exo}[Raison à partir de deux termes]
\explique{Entre deux termes connus, on compte les sauts de $r$ (la différence des rangs),
on trouve $r$, puis on recule jusqu'au rang $0$.}{méthode 3}
\cas{a)}
\begin{calculs}
u_{8} &= u_{3}+5r &&\justif{de $u_{3}$ à $u_{8}$ : $8-3=5$ sauts de $r$}\\
39 &= 14+5r &&\justif{on remplace par les valeurs}\\
5r &= 25 &&\justif{on retire $14$}\\
r &= \res{5} &&\justif{on divise par $5$}
\end{calculs}
\begin{calculs}
u_{0} &= u_{3}-3r &&\justif{on recule de $3$ sauts}\\
u_{0} &= 14-15 = \res{-1} &&\\
u_{n} &= u_{0}+n\,r &&\justif{terme général}\\
u_{n} &= \res{-1+5n} &&
\end{calculs}
\explique{Contrôle : $u_{8}=-1+5\times 8=39$.}{on retrouve la donnée}
\cas{b)}
\begin{calculs}
v_{10} &= v_{2}+8r &&\justif{$10-2=8$ sauts de $r$}\\
6 &= 30+8r &&\\
8r &= -24 &&\justif{on retire $30$}\\
r &= \res{-3} &&
\end{calculs}
\begin{calculs}
v_{0} &= v_{2}-2r &&\justif{on recule de $2$ sauts}\\
v_{0} &= 30+6 = \res{36} &&\justif{$-2\times(-3)=+6$}\\
v_{n} &= \res{36-3n} &&\justif{$v_{n}=v_{0}+n\,r$}
\end{calculs}
\explique{Contrôle : $v_{10}=36-30=6$.}{on retrouve la donnée}
\rem{Reculer quand la raison est négative, c'est remonter : $v_{0}$ est plus grand que
$v_{2}$.}
\cas{c)}
\begin{calculs}
w_{30} &= w_{10}+20r &&\justif{$30-10=20$ sauts de $r$}\\
12 &= 7+20r &&\\
20r &= 5 &&\\
r &= \res{0{,}25} &&\justif{$5\div 20=0{,}25$}
\end{calculs}
\begin{calculs}
w_{0} &= w_{10}-10r &&\justif{on recule de $10$ sauts}\\
w_{0} &= 7-2{,}5 = \res{4{,}5} &&\\
w_{n} &= \res{4{,}5+0{,}25n} &&
\end{calculs}
\explique{Contrôle : $w_{30}=4{,}5+0{,}25\times 30=12$.}{on retrouve la donnée}
\end{exo}

\begin{exo}[Est-ce un terme de la suite ?]
\cas{1.}
\begin{calculs}
u_{n} &= u_{0}+n\,r &&\justif{terme général}\\
u_{n} &= \res{8+3n} &&\justif{$u_{0}=8$ et $r=3$}
\end{calculs}
\cas{2.}
\explique{Un nombre est un terme de la suite s'il existe un \emph{entier naturel} $n$ tel
que $u_{n}$ soit égal à ce nombre.}{le rang $n$ est toujours un entier}
\begin{calculs}
8+3n &= 2\,024 &&\justif{on cherche le rang}\\
3n &= 2\,016 &&\justif{on retire $8$}\\
n &= 672 &&\justif{$2\,016\div 3=672$}
\end{calculs}
\explique{$672$ est un entier : $2\,024=u_{672}$, c'est \res{un terme de la
suite}.}{conclusion}
\begin{calculs}
8+3n &= 2\,025 &&\\
3n &= 2\,017 &&\\
n &= \dfrac{2\,017}{3} &&\justif{$\dfrac{2\,017}{3}\approx 672{,}3$}
\end{calculs}
\explique{$\dfrac{2\,017}{3}$ n'est pas entier : $2\,025$ \res{n'est pas un terme} de la
suite.}{il tombe entre $u_{672}=2\,024$ et $u_{673}=2\,027$}
\cas{3.}
\begin{calculs}
8+3n &> 1\,000 &&\justif{on résout l'inéquation}\\
3n &> 992 &&\\
n &> \dfrac{992}{3} &&\justif{$\dfrac{992}{3}\approx 330{,}7$}
\end{calculs}
\explique{Le premier entier qui convient est $n=331$, et $u_{331}=8+993=1\,001$. Le
premier terme supérieur à $1\,000$ est $\res{u_{331}=1\,001}$.}{$u_{330}=998$ ne convient
pas}
\end{exo}

\begin{exo}[Une cotisation]
\cas{1.}
\explique{$a_{0}=25$ (en 2026) et $a_{n+1}=a_{n}+1{,}5$ : la suite est arithmétique, de
raison $1{,}5$.}{on ajoute $1{,}5$~€ par an}
\begin{calculs}
a_{n} &= a_{0}+n\,r &&\justif{terme général}\\
a_{n} &= \res{25+1{,}5n} &&\\
a_{10} &= 25+1{,}5\times 10 &&\justif{$2036=2026+10$ : $n=10$}\\
a_{10} &= 40 &&
\end{calculs}
\explique{En 2036, la cotisation coûte $\res{40~\text{€}}$.}{on conclut dans le contexte}
\cas{2.}
\begin{calculs}
25+1{,}5n &> 50 &&\justif{la cotisation dépasse $50$~€}\\
1{,}5n &> 25 &&\\
n &> \dfrac{25}{1{,}5} &&\justif{$\dfrac{25}{1{,}5}\approx 16{,}7$}
\end{calculs}
\explique{Le premier entier est $n=17$ : la cotisation dépasse $50$~€ à partir de
$2026+17$, soit \res{en 2043}.}{contrôle : $a_{16}=49$ et $a_{17}=50{,}5$}
\rem{On répond par une année, pas par un rang : on revient toujours au contexte.}
\end{exo}

\partie{Montrer qu'une suite est (ou n'est pas) arithmétique}

\begin{exo}[Avec le terme général]
\explique{On calcule $u_{n+1}-u_{n}$ : si le résultat est un nombre qui ne dépend pas de
$n$, la suite est arithmétique et ce nombre est la raison.}{propriété 2}
\cas{a)}
\begin{calculs}
u_{n+1}-u_{n} &= 5(n+1)+2-(5n+2) &&\justif{$u_{n+1}$ : on remplace $n$ par $n+1$}\\
u_{n+1}-u_{n} &= 5n+5+2-5n-2 &&\justif{on développe}\\
u_{n+1}-u_{n} &= 5 &&
\end{calculs}
\explique{La différence ne dépend pas de $n$ : $(u_{n})$ est arithmétique, de raison
$\res{5}$ et de premier terme $\res{u_{0}=2}$.}{$u_{0}=5\times 0+2$}
\rem{Erreur fréquente : écrire $u_{n+1}=5n+2+1$. On remplace $n$ par $(n+1)$ dans la
formule, on n'ajoute pas $1$ au résultat.}
\cas{b)}
\begin{calculs}
v_{n+1}-v_{n} &= (n+1)^{2}+1-(n^{2}+1) &&\\
v_{n+1}-v_{n} &= n^{2}+2n+1+1-n^{2}-1 &&\justif{identité remarquable $(a+b)^{2}$}\\
v_{n+1}-v_{n} &= 2n+1 &&
\end{calculs}
\explique{La différence dépend de $n$ : $(v_{n})$ \res{n'est pas arithmétique}.}{$v_{1}-v_{0}=1$
mais $v_{2}-v_{1}=3$}
\cas{c)}
\begin{calculs}
w_{n} &= (n+4)^{2}-n^{2} &&\\
w_{n} &= n^{2}+8n+16-n^{2} &&\justif{on développe d'abord}\\
w_{n} &= 8n+16 &&
\end{calculs}
\explique{$w_{n}=8n+16$ est de la forme $an+b$ : $(w_{n})$ est arithmétique, de raison
$\res{8}$ et de premier terme $\res{w_{0}=16}$.}{remarque du cours : $u_{n}=an+b$, raison
$a$}
\rem{Le carré $n^{2}$ disparaît : il ne faut pas conclure trop vite qu'une suite avec
un carré n'est pas arithmétique.}
\cas{d)}
\begin{calculs}
t_{n} &= \dfrac{3-2n}{4} &&\\
t_{n} &= \dfrac{3}{4}-\dfrac{2n}{4} &&\justif{on sépare la fraction}\\
t_{n} &= 0{,}75-0{,}5n &&
\end{calculs}
\explique{De la forme $an+b$ avec $a=-0{,}5$ : $(t_{n})$ est arithmétique, de raison
$\res{-0{,}5}$ et de premier terme $\res{t_{0}=0{,}75}$.}{contrôle : $t_{1}-t_{0}=0{,}25-0{,}75=-0{,}5$}
\end{exo}

\begin{exo}[Avec une relation de récurrence]
\cas{a)}
\explique{$u_{n+1}=u_{n}+7$, donc $u_{n+1}-u_{n}=7$ : la différence est constante.}{la
relation a la forme $u_{n+1}=u_{n}+r$}
\explique{$(u_{n})$ est \res{arithmétique, de raison $7$}.}{premier terme $u_{0}=2$}
\cas{b)}
\explique{La relation multiplie par $2$ avant d'ajouter $1$ : elle n'a pas la forme
$v_{n+1}=v_{n}+r$. On calcule les premiers termes.}{on cherche deux écarts différents}
\begin{calculs}
v_{1} &= 2\times 1+1 = 3 &&\justif{la relation avec $n=0$}\\
v_{2} &= 2\times 3+1 = 7 &&
\end{calculs}
\explique{$v_{1}-v_{0}=2$ mais $v_{2}-v_{1}=4$ : $(v_{n})$ \res{n'est pas
arithmétique}.}{deux écarts différents suffisent}
\cas{c)}
\begin{calculs}
w_{1} &= 3+0 = 3 &&\justif{on ajoute $n=0$}\\
w_{2} &= 3+1 = 4 &&\justif{on ajoute $n=1$}
\end{calculs}
\explique{$w_{1}-w_{0}=0$ mais $w_{2}-w_{1}=1$ : $(w_{n})$ \res{n'est pas
arithmétique}.}{on ajoute $n$, qui change à chaque rang}
\rem{$w_{n+1}=w_{n}+n$ ressemble à la définition, mais ce qu'on ajoute n'est pas une
constante.}
\cas{d)}
\begin{calculs}
t_{n} &= 2(n+3)-n &&\\
t_{n} &= 2n+6-n &&\justif{on développe}\\
t_{n} &= n+6 &&
\end{calculs}
\explique{De la forme $an+b$ avec $a=1$ : $(t_{n})$ est \res{arithmétique, de raison
$1$}.}{ici, une formule explicite, pas une récurrence}
\end{exo}

\begin{exo}[Des carrés d'allumettes]
\cas{1.}
\explique{Pour ajouter un carré à droite, on réutilise le côté commun avec le carré
précédent : il suffit de $3$ allumettes nouvelles. Donc
$a_{n+1}=a_{n}+3$.}{$4\to 7\to 10$ : on ajoute $3$}
\explique{La suite est arithmétique, de raison $3$ ; elle commence à $a_{1}=4$.}{un seul
carré : rang $1$}
\begin{calculs}
a_{n} &= a_{1}+(n-1)\,r &&\justif{le premier terme est $a_{1}$ : $p=1$}\\
a_{n} &= 4+3(n-1) &&\\
a_{n} &= 4+3n-3 &&\justif{on développe}\\
a_{n} &= \res{3n+1} &&
\end{calculs}
\explique{Contrôle : $a_{3}=3\times 3+1=10$.}{on retrouve l'énoncé}
\cas{2.}
\begin{calculs}
a_{50} &= 3\times 50+1 &&\justif{on remplace $n$ par $50$}\\
a_{50} &= 151 &&
\end{calculs}
\explique{Il faut $\res{151}$ allumettes pour $50$ carrés.}{conclusion}
\cas{3.}
\begin{calculs}
3n+1 &\le 200 &&\justif{on ne doit pas dépasser $200$ allumettes}\\
3n &\le 199 &&\\
n &\le \dfrac{199}{3} &&\justif{$\dfrac{199}{3}\approx 66{,}3$}
\end{calculs}
\explique{Le plus grand entier qui convient est $n=66$ : on peut faire au plus
$\res{66~\text{carrés}}$.}{$a_{66}=199$ et $a_{67}=202$}
\end{exo}

\begin{exo}[Somme et produit de suites arithmétiques]
\cas{1.}
\explique{On note $w_{n}=u_{n}+v_{n}$ et on calcule la différence $w_{n+1}-w_{n}$.}{même
critère qu'à l'exercice 9}
\begin{calculs}
w_{n+1}-w_{n} &= (u_{n+1}+v_{n+1})-(u_{n}+v_{n}) &&\justif{on remplace $w$}\\
w_{n+1}-w_{n} &= (u_{n+1}-u_{n})+(v_{n+1}-v_{n}) &&\justif{on regroupe chaque suite}\\
w_{n+1}-w_{n} &= r+r' &&\justif{différences constantes}
\end{calculs}
\explique{$r+r'$ ne dépend pas de $n$ : la suite $(u_{n}+v_{n})$ \res{est arithmétique,
de raison $r+r'$}.}{on additionne les raisons}
\cas{2.}
\explique{$u_{n}\times v_{n}=n\times n=n^{2}$ : les premiers termes sont $0$ ;\; $1$ ;\;
$4$.}{on calcule pour $n=0$, $1$ et $2$}
\explique{Écarts $1$ puis $3$ : la suite $(u_{n}\times v_{n})$ \res{n'est pas
arithmétique}.}{un contre-exemple suffit}
\rem{À retenir : la somme de deux suites arithmétiques est arithmétique, mais pas leur
produit.}
\end{exo}

\partie{Variations : croissance linéaire}

\begin{exo}[Étudier et représenter]
\cas{1.}
\explique{On ajoute $1{,}5$ à chaque terme : $(u_{n})$ est \res{arithmétique, de raison
$r=1{,}5$}.}{définition}
\explique{$r>0$ : la suite est \res{strictement croissante}.}{propriété 3 : le signe de
$r$ donne le sens}
\cas{2. et 3.}
\begin{minipage}[c]{0.58\linewidth}\raggedright
$u_{n}=2+1{,}5n$ : $u_{0}=2$ ;\; $u_{1}=3{,}5$ ;\; $u_{2}=5$ ;\; $u_{3}=6{,}5$ ;\;
$u_{4}=8$.
{\small\itshape\color{black!60}on place les points $(n\,;u_{n})$}

\smallskip
Les points sont alignés sur la droite d'équation $\res{y=1{,}5x+2}$.
{\small\itshape\color{black!60}coefficient directeur $r$, ordonnée à l'origine $u_{0}$ :
la fonction affine associée à $u_{n}=1{,}5n+2$}
\end{minipage}\hfill
\begin{minipage}[c]{0.38\linewidth}\centering
\begin{tikzpicture}[x=0.4cm,y=0.4cm]
  \reperesuite{5}{0}{9}{u_{n}}
  \gradx{1,2,3,4} \grady{2,4,6,8}
  \draw[coulCourbe,dashed] (0,2)--(4.6,8.9);
  \foreach \n in {0,1,...,4}
    \path[coulCourbe] (\n,{2+1.5*\n}) node[croix]{};
\end{tikzpicture}
\end{minipage}
\rem{On ne relie pas les points dans la réponse de l'élève : une suite n'est définie que
pour des entiers. La droite en pointillés montre seulement l'alignement.}
\end{exo}

\begin{exo}[Sens de variation]
\explique{Pour une suite arithmétique, seul le signe de la raison compte : on cherche
donc $r$.}{propriété 3}
\cas{a)}
\explique{$u_{n}=3-4n$ est de la forme $an+b$ : raison $r=-4$.}{le coefficient de $n$}
\explique{$r<0$ : $(u_{n})$ est \res{strictement décroissante}.}{propriété 3}
\cas{b)}
\explique{$v_{n+1}=v_{n}+0{,}2$ : raison $r=0{,}2$ ; $r>0$, $(v_{n})$ est \res{strictement
croissante}.}{propriété 3}
\rem{Le premier terme $v_{0}=-5$ est négatif, mais cela ne change rien : seul le signe
de $r$ compte.}
\cas{c)}
\explique{$w_{n+1}-w_{n}=7-7=0$ : raison $r=0$, $(w_{n})$ est \res{constante}.}{tous les
termes valent $7$}
\cas{d)}
\begin{calculs}
t_{8} &= t_{3}+5r &&\justif{$8-3=5$ sauts de $r$}\\
-5 &= 10+5r &&\\
5r &= -15 &&\\
r &= -3 &&
\end{calculs}
\explique{$r<0$ : $(t_{n})$ est \res{strictement décroissante}.}{on pouvait le voir :
$t_{8}<t_{3}$}
\end{exo}

\begin{exo}[Lire une suite arithmétique]
\cas{1.}
\explique{Le point d'abscisse $0$ a pour ordonnée $1$ : $\res{u_{0}=1}$.}{on lit
l'ordonnée du point d'abscisse $0$}
\explique{D'un point au suivant, on monte de $2$ : $u_{1}=3$, $u_{2}=5$, $u_{3}=7$. La
raison est $\res{r=2}$.}{écart constant entre deux points voisins}
\cas{2.}
\begin{calculs}
u_{n} &= u_{0}+n\,r &&\justif{terme général}\\
u_{n} &= \res{1+2n} &&\\
u_{20} &= 1+2\times 20 &&\\
u_{20} &= \res{41} &&
\end{calculs}
\end{exo}

\begin{exo}[À partir de quel rang ?]
\explique{On résout l'inéquation, puis on prend le premier \emph{entier} qui convient,
car $n$ est un rang.}{méthode 6}
\cas{a)}
\begin{calculs}
4n-9 &> 150 &&\justif{on écrit $u_{n}>150$}\\
4n &> 159 &&\justif{on ajoute $9$}\\
n &> 39{,}75 &&\justif{on divise par $4>0$ : le sens ne change pas}
\end{calculs}
\explique{$u_{n}>150$ à partir du \res{rang $40$}.}{contrôle : $u_{39}=147$ et
$u_{40}=151$}
\cas{b)}
\begin{calculs}
120-7n &< 0 &&\justif{on écrit $v_{n}<0$}\\
-7n &< -120 &&\justif{on retire $120$}\\
n &> \dfrac{120}{7} &&\justif{on divise par $-7<0$ : le sens change}
\end{calculs}
\explique{$\dfrac{120}{7}\approx 17{,}1$ : $v_{n}<0$ à partir du \res{rang $18$}.}{contrôle :
$v_{17}=1$ et $v_{18}=-6$}
\rem{On ne prend pas l'arrondi : $17{,}1$ s'arrondit à $17$, mais $v_{17}=1$ est encore
positif. Il faut le premier entier strictement supérieur.}
\cas{c)}
\begin{calculs}
2{,}5n+1 &\ge 60 &&\\
2{,}5n &\ge 59 &&\\
n &\ge 23{,}6 &&\justif{$59\div 2{,}5=23{,}6$}
\end{calculs}
\explique{$w_{n}\ge 60$ à partir du \res{rang $24$}.}{contrôle : $w_{23}=58{,}5$ et
$w_{24}=61$}
\end{exo}

\partie{Somme de termes d'une suite arithmétique}

\begin{exo}[Sommes d'entiers]
\explique{On utilise $1+2+\dots+n=\dfrac{n(n+1)}{2}$, en lisant $n$ : le dernier
nombre de la somme.}{propriété 4}
\cas{a)}
\begin{calculs}
A &= 1+2+\dots+50 &&\\
A &= \dfrac{50\times 51}{2} &&\justif{$n=50$ dans la formule}\\
A &= \res{1\,275} &&
\end{calculs}
\cas{b)}
\begin{calculs}
B &= 1+2+\dots+150 &&\\
B &= \dfrac{150\times 151}{2} &&\justif{$n=150$}\\
B &= \res{11\,325} &&
\end{calculs}
\cas{c)}
\explique{Tous les termes sont des multiples de $5$ : $5=5\times 1$, $10=5\times 2$\dots\;
$200=5\times 40$.}{la formule ne s'applique qu'à $1+2+\dots+n$}
\begin{calculs}
C &= 5+10+\dots+200 &&\\
C &= 5\times(1+2+\dots+40) &&\justif{on met $5$ en facteur}\\
C &= 5\times\dfrac{40\times 41}{2} &&\justif{$n=40$}\\
C &= 5\times 820 &&\\
C &= \res{4\,100} &&
\end{calculs}
\cas{d)}
\explique{$D$ ne commence pas à $1$ : on calcule $1+2+\dots+200$ et on retire les termes
de $1$ à $100$, qui ne sont pas dans $D$.}{on se ramène à la formule}
\begin{calculs}
D &= (1+\dots+200)-(1+\dots+100) &&\\
D &= \dfrac{200\times 201}{2}-\dfrac{100\times 101}{2} &&\\
D &= 20\,100-5\,050 &&\\
D &= \res{15\,050} &&
\end{calculs}
\rem{On retire jusqu'à $100$, pas jusqu'à $101$ : $101$ fait partie de $D$. Contrôle par
la méthode suivante : $100$ termes, de $101$ à $200$, et
$100\times\dfrac{101+200}{2}=15\,050$.}
\end{exo}

\begin{exo}[Termes consécutifs]
\explique{Pour des termes consécutifs d'une suite arithmétique :
$S=\text{nombre de termes}\times\dfrac{\text{premier}+\text{dernier}}{2}$. Il faut donc
compter les termes avec soin.}{propriété 4}
\cas{a)}
\explique{$u_{n}=4n+1$ est de la forme $an+b$ : la suite est arithmétique, de raison
$4$.}{on peut appliquer la formule}
\explique{De $u_{0}$ à $u_{15}$ : $15+1=16$ termes ; $u_{0}=1$ et
$u_{15}=4\times 15+1=61$.}{nombre de termes, premier et dernier}
\begin{calculs}
S &= 16\times\dfrac{1+61}{2} &&\justif{nombre de termes $\times$ moyenne des extrêmes}\\
S &= 16\times 31 &&\\
S &= \res{496} &&
\end{calculs}
\rem{Erreur fréquente : compter $15$ termes. De $u_{0}$ à $u_{15}$, il y a $u_{0}$ en
plus : $16$ termes.}
\cas{b)}
\explique{$16-9=7$ ;\; $23-16=7$ : suite arithmétique de raison $7$ et de premier terme
$9$, donc $u_{n}=9+7n$.}{on reconnaît la suite}
\begin{calculs}
9+7n &= 128 &&\justif{rang du dernier terme}\\
7n &= 119 &&\\
n &= 17 &&
\end{calculs}
\explique{De $u_{0}$ à $u_{17}$ : $17+1=18$ termes.}{on compte $u_{0}$}
\begin{calculs}
T &= 18\times\dfrac{9+128}{2} &&\\
T &= 9\times 137 &&\justif{$\dfrac{18}{2}=9$}\\
T &= \res{1\,233} &&
\end{calculs}
\cas{c)}
\explique{$47-50=-3$ ;\; $44-47=-3$ : suite arithmétique de raison $-3$, donc
$u_{n}=50-3n$.}{la suite décroît}
\begin{calculs}
50-3n &= 2 &&\justif{rang du dernier terme}\\
-3n &= -48 &&\\
n &= 16 &&
\end{calculs}
\explique{De $u_{0}$ à $u_{16}$ : $17$ termes.}{$16+1$}
\begin{calculs}
U &= 17\times\dfrac{50+2}{2} &&\\
U &= 17\times 26 &&\\
U &= \res{442} &&
\end{calculs}
\end{exo}

\begin{exo}[Somme des nombres impairs]
\cas{1.}
\explique{Les impairs $1$ ;\; $3$ ;\; $5$\dots\ forment une suite arithmétique de raison
$2$ et de premier terme $1$ : $u_{n}=1+2n$.}{on ajoute $2$ pour passer à l'impair
suivant}
\begin{calculs}
1+2n &= 99 &&\justif{rang du dernier terme}\\
2n &= 98 &&\\
n &= 49 &&
\end{calculs}
\explique{De $u_{0}$ à $u_{49}$ : $50$ termes.}{$49+1$}
\begin{calculs}
S &= 50\times\dfrac{1+99}{2} &&\\
S &= 50\times 50 &&\\
S &= \res{2\,500} &&
\end{calculs}
\cas{2.}
\explique{Le dernier terme $2n-1$ s'écrit $1+2(n-1)$ : c'est $u_{n-1}$. De $u_{0}$ à
$u_{n-1}$, il y a $n$ termes.}{$(n-1)+1=n$}
\begin{calculs}
S &= n\times\dfrac{1+(2n-1)}{2} &&\justif{nombre de termes $\times$ moyenne des extrêmes}\\
S &= n\times\dfrac{2n}{2} &&\justif{$1-1=0$}\\
S &= \res{n^{2}} &&
\end{calculs}
\explique{Contrôle : avec $n=50$ (impairs jusqu'à $99$), on retrouve $50^{2}=2\,500$.}{la
question 1}
\end{exo}

\begin{exo}[Les gradins]
\cas{1.}
\explique{On note $u_{n}$ le nombre de places de la rangée $n$ : $u_{1}=20$ et
$u_{n+1}=u_{n}+3$, suite arithmétique de raison $3$.}{$3$ places de plus à chaque
rangée}
\begin{calculs}
u_{n} &= u_{1}+(n-1)\,r &&\justif{le premier terme est $u_{1}$}\\
u_{25} &= 20+24\times 3 &&\justif{de la rangée $1$ à la rangée $25$ : $24$ sauts}\\
u_{25} &= 92 &&
\end{calculs}
\explique{La dernière rangée compte $\res{92}$ places.}{conclusion}
\cas{2.}
\explique{De $u_{1}$ à $u_{25}$ : $25$ termes.}{$25-1+1$}
\begin{calculs}
S &= 25\times\dfrac{20+92}{2} &&\justif{nombre de termes $\times$ moyenne des extrêmes}\\
S &= 25\times 56 &&\\
S &= 1\,400 &&
\end{calculs}
\explique{L'amphithéâtre compte $\res{1\,400}$ places.}{conclusion}
\end{exo}

\partie{Problèmes et suites arithmétiques}

\begin{exo}[Épargner]
\cas{1.}
\begin{calculs}
u_{1} &= 50+10 = \res{60} &&\justif{$10$~€ de plus que le mois précédent}\\
u_{2} &= 60+10 = \res{70} &&
\end{calculs}
\explique{Chaque mois, on ajoute $10$ : $u_{n+1}=u_{n}+10$. La suite est
\res{arithmétique, de raison $10$} et de premier terme $u_{0}=50$.}{on cite la relation}
\cas{2.}
\begin{calculs}
u_{n} &= u_{0}+n\,r &&\justif{terme général}\\
u_{n} &= \res{50+10n} &&
\end{calculs}
\cas{3.}
\begin{calculs}
50+10n &> 200 &&\justif{elle épargne plus de $200$~€}\\
10n &> 150 &&\\
n &> 15 &&
\end{calculs}
\explique{Le premier entier strictement supérieur à $15$ est $\res{n=16}$ : c'est le
17\textsuperscript{e} mois, en mai de l'année suivante.}{janvier : $n=0$ ; janvier
suivant : $n=12$}
\rem{$u_{15}=200$ ne convient pas : \og plus de $200$\fg{} est une inégalité stricte.}
\cas{4.}
\explique{Les $12$ premiers mois : de $u_{0}$ à $u_{11}$, soit $12$ termes ;
$u_{11}=50+110=160$.}{le dernier mois est $n=11$, pas $n=12$}
\begin{calculs}
S &= 12\times\dfrac{50+160}{2} &&\justif{nombre de termes $\times$ moyenne des extrêmes}\\
S &= 12\times 105 &&\\
S &= 1\,260 &&
\end{calculs}
\explique{Sur les $12$ premiers mois, elle aura épargné $\res{1\,260~\text{€}}$.}{on
conclut dans le contexte}
\end{exo}

\begin{exo}[Deux tarifs de vélos en libre-service]
\cas{1.}
\explique{Tarif A : $2$~€ par trajet, donc $\res{a_{n}=2n}$.}{$a_{0}=0$, on ajoute $2$ par
trajet}
\explique{Tarif B : $15$~€ d'abonnement, puis $0{,}80$~€ par trajet, donc
$\res{b_{n}=15+0{,}8n}$.}{$b_{0}=15$, on ajoute $0{,}8$}
\explique{Les deux sont de la forme $an+b$ : ce sont deux \res{suites arithmétiques}, de
raisons $2$ et $0{,}8$.}{remarque du cours}
\cas{2.}
\explique{Le tarif B est le plus avantageux quand il coûte moins cher : $b_{n}<a_{n}$.}{on
traduit la question}
\begin{calculs}
15+0{,}8n &< 2n &&\\
15 &< 1{,}2n &&\justif{on retire $0{,}8n$}\\
12{,}5 &< n &&\justif{$15\div 1{,}2=12{,}5$}
\end{calculs}
\explique{Le tarif B est le plus avantageux \res{à partir de $13$ trajets}.}{contrôle :
$a_{13}=26$ et $b_{13}=25{,}4$}
\end{exo}

\begin{exo}[Un entraînement progressif]
\cas{1.}
\explique{On note $d_{n}$ la distance du $n$-ième jour : $d_{1}=3$ et
$d_{n+1}=d_{n}+0{,}5$, suite arithmétique de raison $0{,}5$.}{la suite commence au rang
$1$}
\begin{calculs}
d_{n} &= d_{1}+(n-1)\,r &&\justif{le premier terme est $d_{1}$}\\
d_{n} &= 3+0{,}5(n-1) &&\\
d_{n} &= 3+0{,}5n-0{,}5 &&\justif{on développe}\\
d_{n} &= \res{0{,}5n+2{,}5} &&
\end{calculs}
\explique{Le $n$-ième jour, il parcourt $0{,}5n+2{,}5$~km.}{contrôle : $d_{1}=3$}
\cas{2.}
\explique{Au bout de $n$ jours : de $d_{1}$ à $d_{n}$, soit $n$ termes ; le premier vaut
$3$ et le dernier $0{,}5n+2{,}5$.}{nombre de termes, premier et dernier}
\begin{calculs}
S_{n} &= n\times\dfrac{3+(0{,}5n+2{,}5)}{2} &&\justif{nombre de termes $\times$ moyenne des
extrêmes}\\
S_{n} &= \res{\dfrac{n(0{,}5n+5{,}5)}{2}} &&\justif{$3+2{,}5=5{,}5$}
\end{calculs}
\cas{3.}
\explique{On cherche le premier $n$ tel que $S_{n}\ge 100$ ; on se ramène à l'inéquation
de l'énoncé.}{on traduit la question}
\begin{calculs}
\dfrac{n(0{,}5n+5{,}5)}{2} &\ge 100 &&\\
n(0{,}5n+5{,}5) &\ge 200 &&\justif{on multiplie par $2$}\\
0{,}5n^{2}+5{,}5n &\ge 200 &&\justif{on développe}\\
n^{2}+11n &\ge 400 &&\justif{on multiplie par $2$}\\
n^{2}+11n-400 &\ge 0 &&\justif{on retire $400$}
\end{calculs}
\explique{On étudie le signe du trinôme $n^{2}+11n-400$ : $a=1$, $b=11$,
$c=-400$.}{chapitre 3 : discriminant}
\begin{calculs}
\Delta &= b^{2}-4ac &&\\
\Delta &= 121+1\,600 &&\justif{$-4\times 1\times(-400)=+1\,600$}\\
\Delta &= 1\,721 &&\justif{$\Delta>0$ : deux racines}
\end{calculs}
\begin{calculs}
n_{1} &= \dfrac{-11-\sqrt{1\,721}}{2} &&\justif{négative : ne convient pas}\\
n_{2} &= \dfrac{-11+\sqrt{1\,721}}{2} &&\justif{$n_{2}\approx 15{,}2$}
\end{calculs}
\explique{$a=1>0$ : le trinôme est positif à l'extérieur des racines. Pour $n\ge 0$, il
est positif quand $n\ge n_{2}\approx 15{,}2$.}{signe de $a$ à l'extérieur des racines}
\explique{Le premier entier est $16$ : il aura parcouru $100$~km \res{au bout de
$16$ jours}.}{contrôle : $S_{15}=97{,}5$ et $S_{16}=108$}
\end{exo}

\begin{exo}[La comète de Halley]
\cas{1.}
\explique{Elle revient tous les $76$ ans : $p_{n+1}=p_{n}+76$, suite arithmétique de
raison $76$ et de premier terme $p_{0}=1\,758$.}{on ajoute $76$ ans}
\begin{calculs}
p_{n} &= p_{0}+n\,r &&\justif{terme général}\\
p_{n} &= \res{1\,758+76n} &&
\end{calculs}
\begin{calculs}
1\,758+76n &= 1\,986 &&\justif{on cherche le rang de 1986}\\
76n &= 228 &&\\
n &= 3 &&\justif{$228\div 76=3$}
\end{calculs}
\explique{$n=3$ est entier : 1986 est bien un passage, le 3\textsuperscript{e} retour
après 1758 ($p_{3}=1\,986$).}{1834, 1910, 1986}
\cas{2.}
\begin{calculs}
p_{4} &= 1\,758+76\times 4 &&\justif{le retour suivant}\\
p_{4} &= 2\,062 &&
\end{calculs}
\explique{Elle repassera \res{en 2062}.}{$1\,986+76=2\,062$}
\cas{3.}
\explique{On remonte le temps : 1066 serait un passage si l'écart $1\,758-1\,066=692$
était un multiple de $76$.}{on prolonge le modèle vers le passé}
\explique{Or $76\times 9=684$ et $76\times 10=760$ : $692$ n'est pas un multiple de $76$.
Selon le modèle, la comète \res{n'a pas pu passer en 1066}.}{$692\div 76\approx 9{,}1$}
\explique{Pourtant elle est bien passée en 1066 : sa période varie en réalité (de $74$ à
$79$ ans). Le modèle n'est qu'une approximation, valable sur quelques retours.}{on
critique le modèle}
\end{exo}

\partie{Suites géométriques}

\begin{exo}[Reconnaître une suite géométrique]
\explique{On divise chaque terme par le précédent : si le quotient est toujours le même,
la suite peut être géométrique et ce quotient est la raison.}{une suite géométrique
multiplie toujours par le même nombre $q$}
\cas{a)}
\explique{$\dfrac{6}{2}=3$ ;\; $\dfrac{18}{6}=3$ ;\; $\dfrac{54}{18}=3$.}{chaque terme
divisé par le précédent}
\explique{Le quotient est toujours $3$ : la suite peut être géométrique, de raison
$\res{q=3}$.}{on multiplie par $3$}
\cas{b)}
\explique{$\dfrac{32}{64}=0{,}5$ ;\; $\dfrac{16}{32}=0{,}5$ ;\; $\dfrac{8}{16}=0{,}5$.}{toujours
le suivant divisé par le précédent}
\explique{La suite peut être géométrique, de raison $\res{q=0{,}5}$.}{multiplier par
$0{,}5$, c'est diviser par $2$}
\rem{Erreur fréquente : répondre $q=2$ en divisant $64$ par $32$. Une suite qui diminue
(à termes positifs) a une raison comprise entre $0$ et $1$.}
\cas{c)}
\explique{$\dfrac{6}{3}=2$ mais $\dfrac{10}{6}=\dfrac{5}{3}$.}{deux quotients différents
suffisent}
\explique{La suite \res{n'est pas géométrique}.}{quotient non constant}
\cas{d)}
\explique{$\dfrac{1\,100}{1\,000}=1{,}1$ ;\; $\dfrac{1\,210}{1\,100}=1{,}1$ ;\;
$\dfrac{1\,331}{1\,210}=1{,}1$.}{chaque terme divisé par le précédent}
\explique{La suite peut être géométrique, de raison $\res{q=1{,}1}$.}{multiplier par
$1{,}1$ : augmenter de $10\,\%$}
\rem{Les écarts ($100$ ; $110$ ; $121$) ne sont pas constants : la suite n'est pas
arithmétique, elle augmente du même pourcentage.}
\end{exo}

\begin{exo}[Traduire une évolution en pourcentage]
\explique{Augmenter de $t\,\%$, c'est multiplier par $1+\dfrac{t}{100}$ ; diminuer de
$t\,\%$, c'est multiplier par $1-\dfrac{t}{100}$ : ce coefficient multiplicateur est la
raison.}{rappel : coefficient multiplicateur}
\cas{a)}
\explique{On note $u_{n}$ le nombre d'habitants au bout de $n$ années : $u_{0}=25\,000$
et $u_{n+1}=1{,}03\,u_{n}$.}{$+3\,\%$ : $1+\dfrac{3}{100}=1{,}03$}
\explique{Suite géométrique de raison $\res{q=1{,}03}$.}{on multiplie toujours par
$1{,}03$}
\rem{Erreur fréquente : $q=0{,}03$ ou $q=3$. Multiplier par $0{,}03$ donnerait $3\,\%$ de
la population, pas la nouvelle population.}
\cas{b)}
\explique{$u_{0}=900$ et $u_{n+1}=0{,}8\,u_{n}$ : raison $\res{q=0{,}8}$.}{$-20\,\%$ :
$1-\dfrac{20}{100}=0{,}8$}
\cas{c)}
\explique{$u_{0}=500$ et $u_{n+1}=2\,u_{n}$ : raison $\res{q=2}$.}{doubler : multiplier
par $2$}
\cas{d)}
\explique{$u_{0}=2\,000\,000$ et $u_{n+1}=0{,}96\,u_{n}$ : raison
$\res{q=0{,}96}$.}{$-4\,\%$ : $1-0{,}04=0{,}96$}
\end{exo}

\begin{exo}[Calculer les premiers termes]
\explique{Dans une suite géométrique, on obtient chaque terme en multipliant le précédent
par la raison : $u_{n+1}=q\times u_{n}$.}{définition}
\cas{a)}
\begin{calculs}
u_{1} &= 2\times 5 = \res{10} &&\justif{$u_{1}=q\times u_{0}$}\\
u_{2} &= 2\times 10 = \res{20} &&\\
u_{3} &= 2\times 20 = \res{40} &&\\
u_{4} &= 2\times 40 = \res{80} &&
\end{calculs}
\cas{b)}
\begin{calculs}
v_{1} &= 0{,}9\times 1\,000 = \res{900} &&\justif{multiplier par $0{,}9$ : baisse de
$10\,\%$}\\
v_{2} &= 0{,}9\times 900 = \res{810} &&\\
v_{3} &= 0{,}9\times 810 = \res{729} &&\\
v_{4} &= 0{,}9\times 729 = \res{656{,}1} &&
\end{calculs}
\rem{La baisse n'est pas toujours de $100$ : $10\,\%$ d'un nombre qui diminue est de plus
en plus petit ($100$, puis $90$, puis $81$\dots).}
\cas{c)}
\explique{Le premier terme est $w_{1}$ : on calcule $w_{2}$ à $w_{5}$.}{on ne commence
pas au rang $0$}
\begin{calculs}
w_{2} &= 1{,}5\times 4 = \res{6} &&\justif{$w_{2}=q\times w_{1}$}\\
w_{3} &= 1{,}5\times 6 = \res{9} &&\\
w_{4} &= 1{,}5\times 9 = \res{13{,}5} &&\\
w_{5} &= 1{,}5\times 13{,}5 = \res{20{,}25} &&
\end{calculs}
\end{exo}

\begin{exo}[Trois termes en progression géométrique]
\explique{Pour trois termes consécutifs d'une suite géométrique, le quotient du deuxième
par le premier est égal au quotient du troisième par le deuxième.}{les deux quotients
valent $q$ ; $x>0$, on peut diviser}
\begin{calculs}
\dfrac{x+6}{x} &= \dfrac{x+18}{x+6} &&\justif{même quotient}\\
(x+6)^{2} &= x(x+18) &&\justif{produits en croix}\\
x^{2}+12x+36 &= x^{2}+18x &&\justif{on développe (identité remarquable à gauche)}\\
36 &= 6x &&\justif{on retire $x^{2}$ et $12x$ des deux côtés}\\
x &= 6 &&\justif{on divise par $6$}
\end{calculs}
\explique{Avec $x=6$, les trois termes sont $\res{6}$ ; $\res{12}$ ; $\res{24}$.}{$x$ ;\;
$x+6$ ;\; $x+18$}
\explique{Vérification : $\dfrac{12}{6}=2$ et $\dfrac{24}{12}=2$. La raison est
$\res{q=2}$.}{on contrôle toujours}
\rem{Les $x^{2}$ s'éliminent : l'équation, qui semblait du second degré, est du premier
degré.}
\end{exo}

\partie{Terme général d'une suite géométrique}

\begin{exo}[Exprimer $u_{n}$ en fonction de $n$]
\explique{Pour aller de $u_{p}$ à $u_{n}$, on multiplie $(n-p)$ fois par la raison :
$u_{n}=u_{p}\times q^{\,n-p}$ ; avec $p=0$, $u_{n}=u_{0}\times q^{n}$.}{propriété 5}
\cas{a)}
\begin{calculs}
u_{n} &= u_{0}\times q^{n} &&\justif{la formule}\\
u_{n} &= \res{2\times 3^{n}} &&\justif{$u_{0}=2$ et $q=3$}\\
u_{6} &= 2\times 3^{6} &&\justif{on remplace $n$ par $6$}\\
u_{6} &= 2\times 729 &&\justif{$3^{6}=729$}\\
u_{6} &= \res{1\,458} &&
\end{calculs}
\rem{Erreur fréquente : $2\times 3^{6}=6^{6}$. La puissance se calcule avant la
multiplication : seul le $3$ est à la puissance $6$.}
\cas{b)}
\explique{$v_{n+1}=2v_{n}$ : suite géométrique de raison $q=2$, de premier terme
$v_{1}=6$.}{on lit la raison dans la relation}
\begin{calculs}
v_{n} &= v_{1}\times q^{\,n-1} &&\justif{le premier terme est $v_{1}$ : $p=1$}\\
v_{n} &= \res{6\times 2^{\,n-1}} &&\\
v_{8} &= 6\times 2^{7} &&\justif{$8-1=7$}\\
v_{8} &= 6\times 128 &&\\
v_{8} &= \res{768} &&
\end{calculs}
\rem{Pourquoi $n-1$ : de $v_{1}$ à $v_{n}$, on multiplie $n-1$ fois par $q$ (de $v_{1}$
à $v_{3}$, deux fois).}
\cas{c)}
\begin{calculs}
w_{n} &= w_{3}\times q^{\,n-3} &&\justif{on part de $w_{3}$ : $p=3$}\\
w_{n} &= \res{40\times 0{,}5^{\,n-3}} &&\\
w_{6} &= 40\times 0{,}5^{3} &&\justif{$6-3=3$}\\
w_{6} &= 40\times 0{,}125 &&\\
w_{6} &= \res{5} &&
\end{calculs}
\explique{Contrôle : $40\to 20\to 10\to 5$, on divise trois fois par $2$.}{multiplier
par $0{,}5$, c'est diviser par $2$}
\end{exo}

\begin{exo}[Raison à partir de deux termes]
\explique{Entre deux termes connus, on compte les multiplications par $q$ (la différence
des rangs), on trouve $q$, puis on recule jusqu'au rang $0$ en divisant par $q$.}{méthode
12}
\cas{a)}
\begin{calculs}
u_{4} &= u_{1}\times q^{3} &&\justif{de $u_{1}$ à $u_{4}$ : $3$ multiplications par $q$}\\
162 &= 6\times q^{3} &&\\
q^{3} &= 27 &&\justif{on divise par $6$}\\
q &= \res{3} &&\justif{le nombre dont le cube vaut $27$ : $3^{3}=27$}
\end{calculs}
\begin{calculs}
u_{0} &= \dfrac{u_{1}}{q} = \dfrac{6}{3} &&\justif{on recule d'un rang : on divise par $q$}\\
u_{0} &= \res{2} &&\\
u_{n} &= \res{2\times 3^{n}} &&\justif{$u_{n}=u_{0}\times q^{n}$}
\end{calculs}
\explique{Contrôle : $u_{4}=2\times 81=162$.}{on retrouve la donnée}
\cas{b)}
\begin{calculs}
v_{4} &= v_{2}\times q^{2} &&\justif{$4-2=2$ multiplications par $q$}\\
16 &= 400\times q^{2} &&\\
q^{2} &= 0{,}04 &&\justif{$16\div 400=0{,}04$}\\
q &= \res{0{,}2} &&\justif{$q>0$ et $0{,}2^{2}=0{,}04$}
\end{calculs}
\begin{calculs}
v_{0} &= \dfrac{v_{2}}{q^{2}} = \dfrac{400}{0{,}04} &&\justif{on recule de $2$ rangs}\\
v_{0} &= \res{10\,000} &&\\
v_{n} &= \res{10\,000\times 0{,}2^{n}} &&
\end{calculs}
\explique{Contrôle : $v_{4}=10\,000\times 0{,}0016=16$.}{$0{,}2^{4}=0{,}0016$}
\rem{Reculer en divisant par $0{,}2$ fait grandir les termes : la suite décroît, donc
$v_{0}$ est le plus grand.}
\cas{c)}
\begin{calculs}
w_{6} &= w_{3}\times q^{3} &&\justif{$6-3=3$ multiplications}\\
40 &= 5\times q^{3} &&\\
q^{3} &= 8 &&\\
q &= \res{2} &&\justif{$2^{3}=8$}
\end{calculs}
\begin{calculs}
w_{0} &= \dfrac{w_{3}}{q^{3}} = \dfrac{5}{8} &&\justif{on recule de $3$ rangs}\\
w_{0} &= \res{0{,}625} &&\\
w_{n} &= \res{0{,}625\times 2^{n}} &&
\end{calculs}
\explique{Contrôle : $w_{3}=0{,}625\times 8=5$.}{on retrouve la donnée}
\end{exo}

\begin{exo}[Intérêts simples ou composés]
\cas{1.}
\explique{Placement S : on ajoute $60$~€ par an, $s_{n+1}=s_{n}+60$. La suite
$(s_{n})$ est arithmétique, de raison $60$ et de premier terme $s_{0}=2\,000$.}{même
somme ajoutée chaque année}
\begin{calculs}
s_{n} &= s_{0}+n\,r &&\\
s_{n} &= \res{2\,000+60n} &&
\end{calculs}
\explique{Placement C : $+3\,\%$ par an, $c_{n+1}=1{,}03\,c_{n}$. La suite $(c_{n})$ est
géométrique, de raison $1{,}03$ et de premier terme $c_{0}=2\,000$.}{même pourcentage
chaque année}
\begin{calculs}
c_{n} &= c_{0}\times q^{n} &&\\
c_{n} &= \res{2\,000\times 1{,}03^{n}} &&
\end{calculs}
\rem{$60$~€, c'est $3\,\%$ de $2\,000$~€ : les intérêts simples sont toujours calculés
sur le capital de départ, les intérêts composés sur le capital de l'année.}
\cas{2.}
\begin{calculs}
s_{10} &= 2\,000+60\times 10 &&\\
s_{10} &= \res{2\,600} &&\\
c_{10} &= 2\,000\times 1{,}03^{10} &&\justif{à la calculatrice}\\
c_{10} &\approx \res{2\,687{,}83} &&\justif{au centime}
\end{calculs}
\explique{Au bout de $10$ ans, le placement C rapporte plus : $87{,}83$~€ de plus que le
placement S.}{$2\,687{,}83-2\,600$}
\end{exo}

\begin{exo}[Plier une feuille]
\cas{1.}
\explique{Chaque pliage double l'épaisseur : $e_{n+1}=2\,e_{n}$, suite géométrique de
raison $2$ et de premier terme $e_{0}=0{,}1$.}{doubler : multiplier par $2$}
\begin{calculs}
e_{n} &= e_{0}\times q^{n} &&\\
e_{n} &= \res{0{,}1\times 2^{n}} &&\\
e_{10} &= 0{,}1\times 2^{10} &&\justif{$2^{10}=1\,024$}\\
e_{10} &= \res{102{,}4} &&
\end{calculs}
\explique{Après $10$ pliages, l'épaisseur est de $102{,}4$~mm, environ $10$~cm.}{conclusion}
\cas{2.}
\explique{On convertit d'abord en mm, l'unité de $e_{n}$ : $8\,849$~m
$=8\,849\,000$~mm.}{$1$~m $=1\,000$~mm}
\explique{À la calculatrice : $e_{26}\approx 6\,710\,886$ et $e_{27}\approx 13\,421\,773$.
Il faudrait $\res{27~\text{pliages}}$ pour dépasser l'Everest.}{$e_{26}<8\,849\,000<e_{27}$}
\explique{Terre-Lune : $384\,400$~km $=384\,400\,000\,000$~mm, soit environ
$3{,}8\times 10^{11}$~mm.}{$1$~km $=1\,000\,000$~mm}
\explique{$e_{41}\approx 2{,}2\times 10^{11}$ et $e_{42}\approx 4{,}4\times 10^{11}$ : il
faudrait $\res{42~\text{pliages}}$.}{$e_{41}<3{,}8\times 10^{11}<e_{42}$}
\rem{Entre l'Everest et la Lune, $15$ pliages seulement : chaque pliage double tout. C'est
la croissance exponentielle.}
\end{exo}

\partie{Montrer qu'une suite est (ou n'est pas) géométrique}

\begin{exo}[Avec le terme général]
\explique{On exprime $u_{n+1}$ et on cherche à faire apparaître $u_{n}$ : si
$u_{n+1}=q\times u_{n}$ avec $q$ constant, la suite est géométrique.}{propriété 6}
\cas{a)}
\begin{calculs}
u_{n+1} &= 5\times 3^{n+1} &&\justif{on exprime le terme suivant}\\
u_{n+1} &= 3\times 5\times 3^{n} &&\justif{$3^{n+1}=3\times 3^{n}$}\\
u_{n+1} &= 3\,u_{n} &&\justif{on reconnaît $u_{n}=5\times 3^{n}$}
\end{calculs}
\explique{$(u_{n})$ est géométrique, de raison $\res{3}$ et de premier terme
$\res{u_{0}=5}$.}{$u_{0}=5\times 3^{0}=5$}
\cas{b)}
\begin{calculs}
v_{n+1} &= 4^{n+2} &&\justif{on remplace $n$ par $n+1$ : $(n+1)+1=n+2$}\\
v_{n+1} &= 4\times 4^{n+1} &&\\
v_{n+1} &= 4\,v_{n} &&
\end{calculs}
\explique{$(v_{n})$ est géométrique, de raison $\res{4}$ et de premier terme
$\res{v_{0}=4}$.}{$v_{0}=4^{1}=4$}
\cas{c)}
\explique{$w_{0}=1$ ;\; $w_{1}=2$ ;\; $w_{2}=5$ : $\dfrac{w_{1}}{w_{0}}=2$ mais
$\dfrac{w_{2}}{w_{1}}=2{,}5$.}{un contre-exemple sur les premiers termes}
\explique{Les quotients diffèrent : $(w_{n})$ \res{n'est pas géométrique}.}{il suffit
d'un contre-exemple}
\cas{d)}
\begin{calculs}
t_{n+1} &= \dfrac{7}{2^{n+1}} &&\justif{on exprime le terme suivant}\\
t_{n+1} &= \dfrac{7}{2\times 2^{n}} &&\justif{$2^{n+1}=2\times 2^{n}$}\\
t_{n+1} &= 0{,}5\times\dfrac{7}{2^{n}} &&\justif{$\dfrac{1}{2}=0{,}5$}\\
t_{n+1} &= 0{,}5\,t_{n} &&
\end{calculs}
\explique{$(t_{n})$ est géométrique, de raison $\res{0{,}5}$ et de premier terme
$\res{t_{0}=7}$.}{$2^{0}=1$}
\end{exo}

\begin{exo}[Avec une relation de récurrence]
\cas{a)}
\explique{$u_{n+1}=1{,}2\,u_{n}$ a la forme $u_{n+1}=q\times u_{n}$ : $(u_{n})$ est
\res{géométrique, de raison $1{,}2$}, de premier terme $u_{0}=3$.}{on multiplie toujours
par $1{,}2$}
\cas{b)}
\explique{$v_{0}=3$ ;\; $v_{1}=4{,}2$ ;\; $v_{2}=5{,}4$ : $\dfrac{4{,}2}{3}=1{,}4$ mais
$\dfrac{5{,}4}{4{,}2}\approx 1{,}29$.}{deux quotients différents}
\explique{$(v_{n})$ \res{n'est pas géométrique} : elle est arithmétique, de raison
$1{,}2$.}{on ajoute $1{,}2$}
\rem{Ne pas confondre $v_{n+1}=v_{n}+1{,}2$ (on ajoute : arithmétique) et
$u_{n+1}=1{,}2\,u_{n}$ (on multiplie : géométrique).}
\cas{c)}
\begin{calculs}
w_{1} &= 2\times 1+1 = 3 &&\justif{la relation avec $n=0$}\\
w_{2} &= 2\times 3+1 = 7 &&
\end{calculs}
\explique{$\dfrac{3}{1}=3$ mais $\dfrac{7}{3}\neq 3$ : $(w_{n})$ \res{n'est pas
géométrique}.}{le $+1$ empêche la forme $q\times w_{n}$}
\cas{d)}
\begin{calculs}
t_{n} &= 3^{n}+3^{n+1} &&\\
t_{n} &= 3^{n}+3\times 3^{n} &&\justif{$3^{n+1}=3\times 3^{n}$}\\
t_{n} &= 4\times 3^{n} &&\justif{$1+3=4$ fois $3^{n}$}
\end{calculs}
\explique{$t_{n}$ est de la forme $t_{0}\times q^{n}$ : $(t_{n})$ est \res{géométrique,
de raison $3$}, de premier terme $t_{0}=4$.}{contrôle : $t_{1}=3+9=12$, soit $3\times t_{0}$}
\end{exo}

\begin{exo}[Arithmétique, géométrique ou aucune des deux ?]
\explique{Arithmétique : la différence $u_{n+1}-u_{n}$ est constante. Géométrique :
$u_{n+1}=q\times u_{n}$. Aucune : deux écarts différents \emph{et} deux quotients
différents.}{les deux critères}
\cas{a)}
\begin{calculs}
u_{n+1}-u_{n} &= 4(n+1)-1-(4n-1) &&\\
u_{n+1}-u_{n} &= 4 &&\justif{$4n+4-1-4n+1$}
\end{calculs}
\explique{La différence est constante : $(u_{n})$ est \res{arithmétique}, de raison
$4$.}{forme $an+b$}
\cas{b)}
\begin{calculs}
v_{n+1} &= 6\times 0{,}5^{n+1} &&\\
v_{n+1} &= 0{,}5\times 6\times 0{,}5^{n} &&\\
v_{n+1} &= 0{,}5\,v_{n} &&
\end{calculs}
\explique{$(v_{n})$ est \res{géométrique}, de raison $0{,}5$.}{forme $v_{0}\times q^{n}$}
\cas{c)}
\explique{$w_{0}=0$ ;\; $w_{1}=1$ ;\; $w_{2}=4$ ;\; $w_{3}=9$. Écarts $1$ puis $3$ : pas
arithmétique.}{différences non constantes}
\explique{$\dfrac{4}{1}=4$ mais $\dfrac{9}{4}=2{,}25$ : pas géométrique. $(w_{n})$ n'est
\res{ni l'une ni l'autre}.}{on ne divise pas par $w_{0}=0$}
\cas{d)}
\explique{$a_{n+1}=a_{n}-2$ : on ajoute toujours $-2$, $(a_{n})$ est
\res{arithmétique}, de raison $-2$.}{forme $a_{n+1}=a_{n}+r$}
\cas{e)}
\explique{$b_{n+1}=5\,b_{n}$ : on multiplie toujours par $5$, $(b_{n})$ est
\res{géométrique}, de raison $5$.}{forme $b_{n+1}=q\times b_{n}$}
\end{exo}

\begin{exo}[Produit et somme de suites géométriques]
\cas{1.}
\explique{On note $w_{n}=u_{n}\times v_{n}$ et on fait apparaître $w_{n}$ dans
$w_{n+1}$.}{propriété 6}
\begin{calculs}
w_{n+1} &= u_{n+1}\times v_{n+1} &&\justif{définition de $w_{n}$ au rang $n+1$}\\
w_{n+1} &= q\,u_{n}\times q'\,v_{n} &&\justif{$u_{n+1}=q\,u_{n}$ et $v_{n+1}=q'\,v_{n}$}\\
w_{n+1} &= qq'\times u_{n}v_{n} &&\justif{on regroupe les raisons}\\
w_{n+1} &= qq'\,w_{n} &&
\end{calculs}
\explique{$(u_{n}\times v_{n})$ \res{est géométrique, de raison $qq'$}.}{on multiplie les
raisons}
\cas{2.}
\explique{$2^{n}+3^{n}$ vaut $2$ ;\; $5$ ;\; $13$ pour $n=0$, $1$ et $2$. Quotients :
$\dfrac{5}{2}=2{,}5$ mais $\dfrac{13}{5}=2{,}6$.}{un contre-exemple}
\explique{La suite $(u_{n}+v_{n})$ \res{n'est pas géométrique}.}{quotients différents}
\rem{À retenir : le produit de deux suites géométriques est géométrique, pas leur
somme. C'est l'inverse des suites arithmétiques (exercice 12).}
\end{exo}

\partie{Variations : croissance exponentielle}

\begin{exo}[Sens de variation]
\explique{On identifie le premier terme et la raison, puis on applique la propriété 7 :
avec $u_{0}>0$, croissante si $q>1$, décroissante si $0<q<1$ ; si $u_{0}<0$, le sens est
inversé.}{propriété 7}
\cas{a)}
\explique{Géométrique de raison $q=1{,}3$ et de premier terme $u_{0}=4$.}{forme
$u_{0}\times q^{n}$}
\explique{$u_{0}>0$ et $q>1$ : $(u_{n})$ est \res{strictement croissante}.}{propriété 7}
\cas{b)}
\explique{Géométrique de raison $q=0{,}6$ et de premier terme $v_{0}=800$.}{forme
$v_{n+1}=q\times v_{n}$}
\explique{$v_{0}>0$ et $0<q<1$ : $(v_{n})$ est \res{strictement décroissante}.}{multiplier
par $0{,}6$ : baisse de $40\,\%$}
\cas{c)}
\explique{Géométrique de raison $q=2$ et de premier terme $w_{0}=-3$.}{forme
$w_{0}\times q^{n}$}
\explique{$q>1$ mais $w_{0}<0$ : le sens est inversé, $(w_{n})$ est \res{strictement
décroissante}.}{$-3$ ;\; $-6$ ;\; $-12$\dots}
\rem{Erreur fréquente : conclure \og croissante\fg{} en regardant seulement $q>1$. Il faut
toujours regarder aussi le signe du premier terme.}
\cas{d)}
\explique{Géométrique de raison $q=0{,}95$ et de premier terme $t_{0}=2$ : $t_{0}>0$ et
$0<q<1$, $(t_{n})$ est \res{strictement décroissante}.}{$0{,}95$ : baisse de $5\,\%$}
\end{exo}

\begin{exo}[Deux villes]
\cas{1.}
\explique{Ville A : on ajoute $500$ habitants par an, suite arithmétique de raison $500$ :
$\res{a_{n}=20\,000+500n}$.}{$a_{n}=a_{0}+n\,r$}
\explique{Ville B : $+3\,\%$ par an, on multiplie par $1{,}03$, suite géométrique de raison
$1{,}03$ : $\res{b_{n}=15\,000\times 1{,}03^{n}}$.}{$b_{n}=b_{0}\times q^{n}$}
\cas{2.}
\explique{On entre les deux suites dans la calculatrice et on lit le tableau autour du
moment où $b_{n}$ rattrape $a_{n}$ :}{application \texttt{Suites}}
\begin{center}
\renewcommand{\arraystretch}{1.3}
\begin{tabular}{|c|c|c|}
\hline
$n$ & $a_{n}$ & $b_{n}$ (arrondi)\\\hline
$27$ & $33\,500$ & $33\,319$\\\hline
$28$ & $34\,000$ & $34\,319$\\\hline
\end{tabular}
\end{center}
\explique{$b_{27}<a_{27}$ mais $b_{28}>a_{28}$ : la ville B dépasse la ville A pour
$n=28$, soit \res{en 2054}.}{$2026+28=2054$}
\rem{La ville B, plus petite au départ, finit par dépasser la ville A : une croissance
exponentielle dépasse tôt ou tard une croissance linéaire.}
\end{exo}

\begin{exo}[Linéaire ou exponentielle ?]
\cas{1.}
\explique{Les croix ($\times$) sont alignées et on monte de $1$ à chaque rang :
$(a_{n})$ est arithmétique, de raison $1$ et de premier terme $a_{0}=1$. Donc
$\res{a_{n}=1+n}$.}{points alignés : croissance linéaire}
\explique{Les $+$ valent $1$ ;\; $2$ ;\; $4$ ;\; $8$ : chaque terme est le double du
précédent. $(g_{n})$ est géométrique, de raison $2$ et de premier terme $g_{0}=1$. Donc
$\res{g_{n}=2^{n}}$.}{quotient constant : $\dfrac{2}{1}=\dfrac{4}{2}=2$}
\cas{2.}
\begin{calculs}
a_{10} &= 1+10 = \res{11} &&\\
g_{10} &= 2^{10} = \res{1\,024} &&
\end{calculs}
\explique{$g_{10}$ est près de cent fois plus grand que $a_{10}$ : la croissance
exponentielle l'emporte largement.}{les deux suites partaient du même point}
\end{exo}

\begin{exo}[La loi de Moore]
\cas{1.}
\explique{Le nombre double à chaque étape (deux ans) : $t_{n+1}=2\,t_{n}$, suite
géométrique de raison $2$ et de premier terme $t_{0}=2\,300$.}{une étape $=$ deux ans}
\begin{calculs}
t_{n} &= t_{0}\times q^{n} &&\\
t_{n} &= \res{2\,300\times 2^{n}} &&\\
t_{25} &= 2\,300\times 2^{25} &&\justif{à la calculatrice}\\
t_{25} &= \res{77\,175\,193\,600} &&\justif{environ $7{,}7\times 10^{10}$}
\end{calculs}
\explique{$n=25$ correspond à $1971+2\times 25$, soit l'année \res{2021}.}{chaque rang
compte deux ans}
\rem{Erreur fréquente : $1971+25=1996$. Ici, $n$ compte des périodes de deux ans.}
\cas{2.}
\explique{À la calculatrice : $t_{18}\approx 6{,}0\times 10^{8}$, inférieur à un milliard,
et $t_{19}\approx 1{,}2\times 10^{9}$, supérieur.}{un milliard $=10^{9}$}
\explique{On dépasse le milliard pour $n=19$, soit $1971+38$ : \res{en 2009}.}{$2\times 19=38$}
\cas{3.}
\explique{Le nombre est multiplié par le même facteur ($2$) à chaque étape : c'est une
suite géométrique de raison $2>1$, dont les termes augmentent de plus en plus vite.
C'est une croissance exponentielle.}{même pourcentage ($+100\,\%$) à chaque étape}
\end{exo}

\partie{Somme de termes d'une suite géométrique}

\begin{exo}[Sommes de puissances]
\explique{On utilise $1+q+\dots+q^{n}=\dfrac{1-q^{n+1}}{1-q}$, ou plus généralement :
premier terme $\times\dfrac{1-q^{\,\text{nombre de termes}}}{1-q}$.}{propriété 8}
\cas{a)}
\explique{$q=2$ et $n=11$ : de $2^{0}$ à $2^{11}$, il y a $12$ termes.}{l'exposant du
haut est le nombre de termes}
\begin{calculs}
A &= 1+2+2^{2}+\dots+2^{11} &&\\
A &= \dfrac{1-2^{12}}{1-2} &&\justif{exposant $n+1=12$}\\
A &= \dfrac{1-4\,096}{-1} &&\justif{$2^{12}=4\,096$}\\
A &= \res{4\,095} &&
\end{calculs}
\cas{b)}
\explique{$B$ ne commence pas par $1$ : premier terme $3$, raison $3$, et $8$ termes
(exposants de $1$ à $8$).}{on met le premier terme en facteur}
\begin{calculs}
B &= 3+3^{2}+\dots+3^{8} &&\\
B &= 3\times(1+3+\dots+3^{7}) &&\justif{premier terme $3$ en facteur}\\
B &= 3\times\dfrac{1-3^{8}}{1-3} &&\justif{$8$ termes, de $3^{0}$ à $3^{7}$}\\
B &= 3\times\dfrac{1-6\,561}{-2} &&\justif{$3^{8}=6\,561$}\\
B &= 3\times 3\,280 &&\\
B &= \res{9\,840} &&
\end{calculs}
\rem{Erreur fréquente : écrire $B=\dfrac{1-3^{9}}{1-3}$, qui compte le terme $1$ absent
de $B$ (on trouverait $9\,841$).}
\cas{c)}
\explique{$q=0{,}2$ ; de $0{,}2^{0}$ à $0{,}2^{5}$ : $6$ termes.}{exposant $5+1=6$}
\begin{calculs}
C &= 1+0{,}2+\dots+0{,}2^{5} &&\\
C &= \dfrac{1-0{,}2^{6}}{1-0{,}2} &&\\
C &= \dfrac{0{,}999\,936}{0{,}8} &&\justif{$0{,}2^{6}=0{,}000\,064$}\\
C &\approx \res{1{,}250} &&\justif{au millième ($1{,}249\,92$)}
\end{calculs}
\end{exo}

\begin{exo}[Somme de termes consécutifs]
\cas{a)}
\explique{$(u_{n})$ est géométrique, de raison $2$ et de premier terme $u_{0}=3$. De
$u_{0}$ à $u_{7}$ : $8$ termes.}{on compte $u_{0}$}
\begin{calculs}
S &= u_{0}+u_{1}+\dots+u_{7} &&\\
S &= 3\times\dfrac{1-2^{8}}{1-2} &&\justif{premier terme $\times\dfrac{1-q^{8}}{1-q}$}\\
S &= 3\times\dfrac{1-256}{-1} &&\\
S &= 3\times 255 &&\\
S &= \res{765} &&
\end{calculs}
\cas{b)}
\explique{$(v_{n})$ est géométrique, de raison $0{,}5$ et de premier terme
$v_{0}=1\,000$. De $v_{0}$ à $v_{5}$ : $6$ termes.}{$5+1=6$}
\begin{calculs}
T &= v_{0}+v_{1}+\dots+v_{5} &&\\
T &= 1\,000\times\dfrac{1-0{,}5^{6}}{1-0{,}5} &&\\
T &= 1\,000\times\dfrac{0{,}984\,375}{0{,}5} &&\justif{$0{,}5^{6}=0{,}015\,625$}\\
T &= \res{1\,968{,}75} &&\justif{diviser par $0{,}5$, c'est multiplier par $2$}
\end{calculs}
\explique{Contrôle : $1\,000+500+250+125+62{,}5+31{,}25=1\,968{,}75$.}{six termes, on peut
les additionner}
\end{exo}

\begin{exo}[Les grains de riz]
\cas{1.}
\explique{On note $u_{k}$ le nombre de grains sur la case $k$ : $u_{1}=1$ et on double à
chaque case, suite géométrique de raison $2$.}{la suite commence au rang $1$}
\begin{calculs}
u_{k} &= u_{1}\times q^{\,k-1} &&\justif{le premier terme est $u_{1}$}\\
u_{k} &= 2^{\,k-1} &&\\
u_{10} &= 2^{9} &&\\
u_{10} &= 512 &&
\end{calculs}
\explique{La 10\textsuperscript{e} case porte $\res{512}$ grains.}{$9$ doublements depuis la
première case}
\cas{2.}
\explique{Sur les $10$ premières cases : $1+2+\dots+2^{9}$, soit $10$ termes.}{de $2^{0}$
à $2^{9}$}
\begin{calculs}
S &= \dfrac{1-2^{10}}{1-2} &&\justif{propriété 8}\\
S &= \dfrac{1-1\,024}{-1} &&\\
S &= \res{1\,023} &&
\end{calculs}
\explique{Sur les $64$ cases : $1+2+\dots+2^{63}$, soit $64$ termes.}{la case $64$ porte
$2^{63}$ grains}
\begin{calculs}
S' &= \dfrac{1-2^{64}}{1-2} &&\\
S' &= \dfrac{1-2^{64}}{-1} &&\\
S' &= \res{2^{64}-1} &&\justif{environ $1{,}8\times 10^{19}$ grains}
\end{calculs}
\rem{Diviser par $-1$ change tous les signes : $\dfrac{1-2^{64}}{-1}=2^{64}-1$.}
\end{exo}

\begin{exo}[L'escargot]
\cas{1.}
\explique{Les distances parcourues chaque heure sont $1$ ;\; $0{,}5$ ;\; $0{,}25$\dots\ :
suite géométrique de raison $0{,}5$ et de premier terme $1$.}{la moitié de l'heure
précédente : on multiplie par $0{,}5$}
\explique{Au bout de $n$ heures, on additionne $n$ termes : $1+0{,}5+\dots+0{,}5^{\,n-1}$.}{de
$0{,}5^{0}$ à $0{,}5^{\,n-1}$ : $n$ termes}
\begin{calculs}
S_{n} &= 1+0{,}5+\dots+0{,}5^{\,n-1} &&\\
S_{n} &= \dfrac{1-0{,}5^{n}}{1-0{,}5} &&\justif{exposant : le nombre de termes $n$}\\
S_{n} &= \dfrac{1-0{,}5^{n}}{0{,}5} &&\\
S_{n} &= 2\left(1-0{,}5^{n}\right) &&\justif{diviser par $0{,}5$, c'est multiplier par $2$}
\end{calculs}
\cas{2.}
\begin{calculs}
S_{10} &= 2\left(1-0{,}5^{10}\right) &&\\
S_{10} &= 2\left(1-\dfrac{1}{1\,024}\right) &&\justif{$0{,}5^{10}=\dfrac{1}{2^{10}}$}\\
S_{10} &\approx \res{1{,}998} &&
\end{calculs}
\explique{$0{,}5^{n}>0$, donc $1-0{,}5^{n}<1$ et $S_{n}<2$ pour tout $n$ : l'escargot
\res{n'atteindra jamais $2$~m}.}{il manque toujours $2\times 0{,}5^{n}$}
\cas{3.}
\explique{$0{,}5^{n}$ devient aussi proche de $0$ que l'on veut, donc $S_{n}$ se rapproche
de $2$ : on conjecture $\res{\lim\limits_{n\to+\infty}S_{n}=2}$.}{par exemple
$0{,}5^{20}<0{,}000\,001$}
\end{exo}

\partie{Problèmes et suites géométriques}

\begin{exo}[Les rebonds d'une balle]
\cas{1.}
\explique{Remonter à $70\,\%$ de la hauteur, c'est multiplier par $0{,}7$ :
$h_{n+1}=0{,}7\,h_{n}$. La suite est \res{géométrique, de raison $0{,}7$} et de premier
terme $h_{0}=2$.}{coefficient multiplicateur $0{,}7$}
\begin{calculs}
h_{n} &= h_{0}\times q^{n} &&\justif{terme général}\\
h_{n} &= \res{2\times 0{,}7^{n}} &&
\end{calculs}
\cas{2.}
\begin{calculs}
h_{5} &= 2\times 0{,}7^{5} &&\justif{après $5$ rebonds : $n=5$}\\
h_{5} &= 2\times 0{,}168\,07 &&\\
h_{5} &\approx 0{,}34 &&\justif{au cm : $0{,}336\,14$~m}
\end{calculs}
\explique{Après $5$ rebonds, la balle remonte à environ $\res{0{,}34~\text{m}}$, soit
$34$~cm.}{conclusion}
\cas{3.}
\begin{minipage}[t]{0.42\linewidth}\vspace{0pt}
\begin{python}
def seuil():
    h = 2
    n = 0
    while h >= 0.1 :
        h = 0.7*h
        n = n + 1
    return n
\end{python}
\end{minipage}\hfill
\begin{minipage}[t]{0.54\linewidth}\vspace{0pt}\raggedright
{\small\itshape\color{black!60}$10$~cm $=0{,}1$~m : tant que la hauteur n'est pas
inférieure à $0{,}1$, on passe au rebond suivant (on multiplie par $0{,}7$ et on
compte)}

\smallskip
$h_{8}\approx 0{,}115$ et $h_{9}\approx 0{,}081$ : \texttt{seuil()} renvoie $9$, la
hauteur passe sous $10$~cm \res{après le 9\textsuperscript{e} rebond}.
\end{minipage}
\rem{La condition du \texttt{while} est le contraire de ce qu'on cherche : on cherche
$h<0{,}1$, on boucle tant que $h\ge 0{,}1$.}
\end{exo}

\begin{exo}[Doubler un placement]
\cas{1.}
\explique{$+2{,}5\,\%$ par an : on multiplie par $1{,}025$. Suite géométrique de raison
$1{,}025$ et de premier terme $c_{0}=5\,000$.}{$1+\dfrac{2{,}5}{100}=1{,}025$}
\begin{calculs}
c_{n} &= c_{0}\times q^{n} &&\\
c_{n} &= \res{5\,000\times 1{,}025^{n}} &&\\
c_{10} &= 5\,000\times 1{,}025^{10} &&\justif{à la calculatrice}\\
c_{10} &\approx \res{6\,400{,}42} &&\justif{au centime}
\end{calculs}
\cas{2.}
\explique{Le capital a doublé quand $c_{n}\ge 10\,000$. À la calculatrice :
$c_{28}\approx 9\,982{,}48$ et $c_{29}\approx 10\,232{,}04$.}{tableau de valeurs}
\explique{Le capital aura doublé \res{au bout de $29$ ans}.}{premier rang où
$c_{n}\ge 10\,000$}
\rem{Le temps de doublement ne dépend pas du capital : on cherche toujours
$1{,}025^{n}\ge 2$.}
\end{exo}

\begin{exo}[Téléchargements d'une application]
\cas{1.}
\explique{$+8\,\%$ par mois : $u_{n+1}=1{,}08\,u_{n}$, suite géométrique de raison $1{,}08$
et de premier terme $u_{0}=2\,000$.}{coefficient multiplicateur $1{,}08$}
\begin{calculs}
u_{n} &= u_{0}\times q^{n} &&\\
u_{n} &= \res{2\,000\times 1{,}08^{n}} &&\\
u_{11} &= 2\,000\times 1{,}08^{11} &&\justif{à la calculatrice}\\
u_{11} &\approx \res{4\,663} &&\justif{à l'unité}
\end{calculs}
\cas{2.}
\explique{Les $12$ premiers mois : de $u_{0}$ à $u_{11}$, soit $12$ termes.}{le dernier
mois est $n=11$}
\begin{calculs}
S &= 2\,000\times\dfrac{1-1{,}08^{12}}{1-1{,}08} &&\justif{premier terme
$\times\dfrac{1-q^{12}}{1-q}$}\\
S &\approx 37\,954 &&\justif{à la calculatrice}
\end{calculs}
\explique{Sur les $12$ premiers mois, l'application est téléchargée environ
$\res{37\,954}$ fois.}{on conclut dans le contexte}
\rem{$1-1{,}08$ est négatif, comme $1-1{,}08^{12}$ : le quotient est bien positif. Bien
mettre les parenthèses à la calculatrice.}
\end{exo}

\begin{exo}[Deux offres d'emploi]
\explique{On note $a_{n}$ et $b_{n}$ les salaires de l'année $n+1$ : la première année
correspond à $n=0$, la dixième à $n=9$.}{on nomme les suites}
\explique{Offre A : $+1\,000$~€ par an, $a_{n}=24\,000+1\,000n$ (arithmétique). Offre B :
$+6\,\%$ par an, $b_{n}=22\,000\times 1{,}06^{n}$ (géométrique).}{on ajoute ; on multiplie
par $1{,}06$}
\cas{1.}
\explique{Première année : $a_{0}=24\,000$ et $b_{0}=22\,000$ ; c'est l'offre \res{A} qui
rapporte le plus.}{$24\,000>22\,000$}
\begin{calculs}
a_{9} &= 24\,000+1\,000\times 9 = 33\,000 &&\justif{dixième année : $n=9$}\\
b_{9} &= 22\,000\times 1{,}06^{9} &&\justif{à la calculatrice}\\
b_{9} &\approx 37\,169 &&
\end{calculs}
\explique{La dixième année, c'est l'offre \res{B} qui rapporte le plus.}{$37\,169>33\,000$}
\cas{2.}
\explique{Sur dix ans : de $a_{0}$ à $a_{9}$ (et de $b_{0}$ à $b_{9}$), soit $10$
termes.}{$9+1=10$}
\begin{calculs}
S_{A} &= 10\times\dfrac{24\,000+33\,000}{2} &&\justif{arithmétique : nombre de termes
$\times$ moyenne des extrêmes}\\
S_{A} &= 10\times 28\,500 &&\\
S_{A} &= \res{285\,000} &&
\end{calculs}
\begin{calculs}
S_{B} &= 22\,000\times\dfrac{1-1{,}06^{10}}{1-1{,}06} &&\justif{géométrique : premier
terme $\times\dfrac{1-q^{10}}{1-q}$}\\
S_{B} &\approx \res{289\,977} &&\justif{à la calculatrice}
\end{calculs}
\explique{Sur dix ans, l'offre B rapporte environ $4\,977$~€ de plus que l'offre A, bien
qu'elle commence plus bas.}{conclusion : la croissance exponentielle finit par
l'emporter}
\rem{Deux formules de somme différentes : on reconnaît d'abord la nature de chaque
suite.}
\end{exo}

\partie{Suites arithmético-géométriques}

\begin{exo}[Avec une suite auxiliaire]
\cas{1.}
\begin{calculs}
u_{1} &= 0{,}8\times 100+30 = \res{110} &&\justif{la relation avec $n=0$}\\
u_{2} &= 0{,}8\times 110+30 = \res{118} &&\justif{$88+30$}
\end{calculs}
\explique{$u_{1}-u_{0}=10$ mais $u_{2}-u_{1}=8$ : pas arithmétique.}{deux écarts
différents}
\explique{$\dfrac{u_{1}}{u_{0}}=1{,}1$ mais $\dfrac{u_{2}}{u_{1}}\approx 1{,}07$ : pas
géométrique.}{deux quotients différents}
\cas{2.}
\explique{On part de $v_{n+1}$ et on cherche à faire apparaître $v_{n}$.}{même démarche
que la méthode 18}
\begin{calculs}
v_{n+1} &= u_{n+1}-150 &&\justif{définition de $v_{n}$, au rang $n+1$}\\
v_{n+1} &= 0{,}8\,u_{n}+30-150 &&\justif{relation de récurrence de $(u_{n})$}\\
v_{n+1} &= 0{,}8\,u_{n}-120 &&\\
v_{n+1} &= 0{,}8(u_{n}-150) &&\justif{on factorise par $0{,}8$ : $0{,}8\times 150=120$}\\
v_{n+1} &= 0{,}8\,v_{n} &&\justif{on reconnaît $v_{n}=u_{n}-150$}
\end{calculs}
\begin{calculs}
v_{0} &= u_{0}-150 &&\justif{premier terme}\\
v_{0} &= 100-150 = \res{-50} &&
\end{calculs}
\explique{$(v_{n})$ est \res{géométrique, de raison $0{,}8$} et de premier terme
$v_{0}=-50$.}{forme $v_{n+1}=q\times v_{n}$}
\rem{Erreur fréquente : remplacer $u_{n+1}$ par $v_{n+1}+150$ et tourner en rond. On
remplace $u_{n+1}$ par la relation de l'énoncé.}
\cas{3.}
\begin{calculs}
v_{n} &= v_{0}\times q^{n} &&\\
v_{n} &= \res{-50\times 0{,}8^{n}} &&\\
u_{n} &= v_{n}+150 &&\justif{car $v_{n}=u_{n}-150$}\\
u_{n} &= \res{150-50\times 0{,}8^{n}} &&
\end{calculs}
\explique{Contrôle : $u_{2}=150-50\times 0{,}64=118$.}{on retrouve la question 1}
\cas{4.}
\explique{$(v_{n})$ est géométrique avec $v_{0}<0$ et $0<q<1$ : le sens est inversé, elle
est strictement croissante.}{propriété 7 : $-50$ ;\; $-40$ ;\; $-32$\dots}
\explique{$u_{n}=v_{n}+150$ : $(u_{n})$ est aussi \res{strictement
croissante}.}{ajouter $150$ ne change pas le sens}
\explique{$0{,}8^{n}$ devient aussi proche de $0$ que l'on veut, donc $v_{n}$ se rapproche
de $0$ et $u_{n}$ de $150$ : on conjecture
$\res{\lim\limits_{n\to+\infty}u_{n}=150}$.}{$u_{n}$ reste sous $150$}
\end{exo}

\begin{exo}[Une suite qui s'envole]
\cas{1.}
\begin{calculs}
v_{n+1} &= u_{n+1}-2 &&\justif{définition de $v_{n}$, au rang $n+1$}\\
v_{n+1} &= 3u_{n}-4-2 &&\justif{relation de récurrence}\\
v_{n+1} &= 3u_{n}-6 &&\\
v_{n+1} &= 3(u_{n}-2) &&\justif{on factorise par $3$ : $3\times 2=6$}\\
v_{n+1} &= 3\,v_{n} &&
\end{calculs}
\begin{calculs}
v_{0} &= u_{0}-2 = \res{2} &&\justif{$4-2$}
\end{calculs}
\explique{$(v_{n})$ est \res{géométrique, de raison $3$} et de premier terme
$v_{0}=2$.}{forme $v_{n+1}=q\times v_{n}$}
\cas{2.}
\begin{calculs}
v_{n} &= 2\times 3^{n} &&\justif{$v_{0}\times q^{n}$}\\
u_{n} &= v_{n}+2 &&\justif{car $v_{n}=u_{n}-2$}\\
u_{n} &= \res{2+2\times 3^{n}} &&
\end{calculs}
\explique{Contrôle : $u_{1}=3\times 4-4=8$ et $2+2\times 3=8$.}{on calcule $u_{1}$ des
deux façons}
\cas{3.}
\explique{$v_{0}>0$ et $q=3>1$ : $(v_{n})$ est strictement croissante, donc $(u_{n})$
aussi : \res{strictement croissante}.}{propriété 7, puis on ajoute $2$}
\explique{$3^{n}$ devient aussi grand que l'on veut : on conjecture
$\res{\lim\limits_{n\to+\infty}u_{n}=+\infty}$.}{la suite s'envole}
\end{exo}

\begin{exo}[Les poissons d'un étang]
\cas{1.}
\explique{Si $30\,\%$ des poissons disparaissent, il en reste $70\,\%$ : on multiplie par
$0{,}7$. Puis on en introduit $450$ : on ajoute $450$. Donc
$u_{n+1}=0{,}7\,u_{n}+450$.}{d'abord on multiplie, puis on ajoute}
\cas{2.}
\begin{calculs}
v_{n+1} &= u_{n+1}-1\,500 &&\justif{définition de $v_{n}$, au rang $n+1$}\\
v_{n+1} &= 0{,}7\,u_{n}+450-1\,500 &&\justif{relation de récurrence}\\
v_{n+1} &= 0{,}7\,u_{n}-1\,050 &&\\
v_{n+1} &= 0{,}7(u_{n}-1\,500) &&\justif{$0{,}7\times 1\,500=1\,050$}\\
v_{n+1} &= 0{,}7\,v_{n} &&
\end{calculs}
\explique{$(v_{n})$ est géométrique, de raison $0{,}7$ et de premier terme
$v_{0}=2\,000-1\,500=500$.}{$u_{0}=2\,000$}
\begin{calculs}
v_{n} &= 500\times 0{,}7^{n} &&\justif{$v_{0}\times q^{n}$}\\
u_{n} &= v_{n}+1\,500 &&\justif{car $v_{n}=u_{n}-1\,500$}\\
u_{n} &= \res{1\,500+500\times 0{,}7^{n}} &&\justif{on retrouve l'expression demandée}
\end{calculs}
\cas{3.}
\explique{$v_{0}>0$ et $0<q<1$ : $(v_{n})$, donc $(u_{n})$, est \res{strictement
décroissante}.}{propriété 7, puis on ajoute $1\,500$}
\explique{$0{,}7^{n}$ se rapproche de $0$ : on conjecture
$\res{\lim\limits_{n\to+\infty}u_{n}=1\,500}$.}{la population se stabilise}
\cas{4.}
\begin{calculs}
1\,500+500\times 0{,}7^{n} &< 1\,550 &&\justif{moins de $1\,550$ poissons}\\
500\times 0{,}7^{n} &< 50 &&\justif{on retire $1\,500$}\\
0{,}7^{n} &< 0{,}1 &&\justif{on divise par $500$}
\end{calculs}
\explique{À la calculatrice : $0{,}7^{6}\approx 0{,}118$ et $0{,}7^{7}\approx 0{,}082$.
Il y aura moins de $1\,550$ poissons \res{au bout de $7$ ans}.}{contrôle :
$u_{6}\approx 1\,559$ et $u_{7}\approx 1\,541$}
\end{exo}

\begin{exo}[Trouver la suite auxiliaire]
\cas{1.}
\begin{calculs}
\ell &= 0{,}6\,\ell+4 &&\\
\ell-0{,}6\,\ell &= 4 &&\justif{on regroupe les $\ell$}\\
0{,}4\,\ell &= 4 &&\\
\ell &= \res{10} &&\justif{$4\div 0{,}4=10$}
\end{calculs}
\cas{2.}
\begin{calculs}
v_{n+1} &= u_{n+1}-10 &&\justif{$v_{n}=u_{n}-10$, au rang $n+1$}\\
v_{n+1} &= 0{,}6\,u_{n}+4-10 &&\justif{relation de récurrence}\\
v_{n+1} &= 0{,}6\,u_{n}-6 &&\\
v_{n+1} &= 0{,}6(u_{n}-10) &&\justif{$0{,}6\times 10=6$}\\
v_{n+1} &= 0{,}6\,v_{n} &&
\end{calculs}
\explique{$(v_{n})$ est \res{géométrique, de raison $0{,}6$} et de premier terme
$v_{0}=2-10=-8$.}{forme $v_{n+1}=q\times v_{n}$}
\cas{3.}
\begin{calculs}
v_{n} &= -8\times 0{,}6^{n} &&\justif{$v_{0}\times q^{n}$}\\
u_{n} &= v_{n}+10 &&\\
u_{n} &= \res{10-8\times 0{,}6^{n}} &&
\end{calculs}
\explique{$0{,}6^{n}$ se rapproche de $0$ : on conjecture
$\res{\lim\limits_{n\to+\infty}u_{n}=10}$.}{contrôle : $u_{1}=0{,}6\times 2+4=5{,}2$ et
$10-8\times 0{,}6=5{,}2$}
\rem{Pourquoi ce $\ell$ : c'est la valeur que la relation ne change pas (si $u_{n}=10$,
alors $u_{n+1}=0{,}6\times 10+4=10$). C'est grâce à lui que la factorisation tombe juste
à la question 2, et c'est aussi la limite.}
\end{exo}

\partiesansnum{Exercices type bac}

\begin{exo}[Automatismes (QCM)]
\cas{1.}
\explique{$u_{n}=7-2n$, donc $u_{10}=7-20=-13$ : réponse \res{a)}.}{$u_{n}=u_{0}+n\,r$}
\rem{La réponse d) ($-20$) oublie le premier terme $7$.}
\cas{2.}
\explique{Hausse de $4\,\%$ : on multiplie par $1+\dfrac{4}{100}=1{,}04$ : réponse
\res{d)}.}{coefficient multiplicateur}
\cas{3.}
\explique{$1+2+\dots+40=\dfrac{40\times 41}{2}=820$ : réponse \res{b)}.}{$n=40$ ; la
réponse d) oublie de diviser par $2$}
\cas{4.}
\explique{Suite géométrique de raison $2$ et de premier terme $3$ :
$v_{n}=3\times 2^{n}$, réponse \res{c)}.}{$v_{n}=v_{0}\times q^{n}$}
\rem{La réponse b) échange le premier terme et la raison ; d) multiplie $3$ et $2$ avant
la puissance.}
\cas{5.}
\explique{De $3^{0}$ à $3^{5}$ : $6$ termes, donc
$\dfrac{1-3^{6}}{1-3}$ : réponse \res{a)}.}{l'exposant du haut est le nombre de termes}
\end{exo}

\begin{exo}[Un apiculteur]
\textbf{Partie A.}
\cas{1.}
\explique{Il installe $5$ ruches de plus chaque année : $r_{n+1}=r_{n}+5$. La suite
$(r_{n})$ est \res{arithmétique, de raison $5$} et de premier terme
$r_{0}=12$.}{on ajoute toujours $5$}
\begin{calculs}
r_{n} &= r_{0}+5n &&\justif{terme général}\\
r_{n} &= \res{12+5n} &&
\end{calculs}
\cas{2.}
\begin{calculs}
12+5n &\ge 60 &&\justif{au moins $60$ ruches}\\
5n &\ge 48 &&\\
n &\ge 9{,}6 &&
\end{calculs}
\explique{Le premier entier est $n=10$ : il aura au moins $60$ ruches \res{en 2036}.}{contrôle :
$r_{9}=57$ et $r_{10}=62$}
\textbf{Partie B.}
\cas{3.}
\explique{$-4\,\%$ par an : on multiplie par $0{,}96$, $m_{n+1}=0{,}96\,m_{n}$. La suite
$(m_{n})$ est \res{géométrique, de raison $0{,}96$} et de premier terme
$m_{0}=25$.}{$1-\dfrac{4}{100}=0{,}96$}
\begin{calculs}
m_{n} &= m_{0}\times q^{n} &&\\
m_{n} &= \res{25\times 0{,}96^{n}} &&
\end{calculs}
\cas{4.}
\begin{calculs}
m_{10} &= 25\times 0{,}96^{10} &&\justif{$2036=2026+10$}\\
m_{10} &\approx 16{,}6 &&\justif{au dixième}
\end{calculs}
\explique{En 2036, une ruche produit environ $\res{16{,}6~\text{kg}}$ de miel.}{conclusion}
\cas{5.}
\begin{minipage}[t]{0.42\linewidth}\vspace{0pt}
\begin{python}
def seuil():
    m = 25
    n = 0
    while m >= 20 :
        m = 0.96*m
        n = n + 1
    return n
\end{python}
\end{minipage}\hfill
\begin{minipage}[t]{0.54\linewidth}\vspace{0pt}\raggedright
{\small\itshape\color{black!60}tant que la production n'est pas passée sous $20$~kg, on
passe à l'année suivante (on multiplie par $0{,}96$ et on compte)}

\smallskip
$m_{5}\approx 20{,}4$ et $m_{6}\approx 19{,}6$ : \texttt{seuil()} renvoie $6$, la
production passe sous $20$~kg \res{en 2032}.
\end{minipage}
\cas{6.}
\explique{De 2026 à 2035 : de $m_{0}$ à $m_{9}$, soit $10$ termes.}{on compte les
années}
\begin{calculs}
S &= 25\times\dfrac{1-0{,}96^{10}}{1-0{,}96} &&\justif{premier terme
$\times\dfrac{1-q^{10}}{1-q}$}\\
S &\approx 209{,}5 &&\justif{à la calculatrice}
\end{calculs}
\explique{De 2026 à 2035, une ruche produit en tout environ $\res{209{,}5~\text{kg}}$ de
miel.}{conclusion}
\end{exo}

\begin{exo}[Trottinettes en libre-service]
\cas{1.}
\begin{calculs}
u_{1} &= 0{,}95\times 800+60 &&\justif{on retire $5\,\%$, puis on ajoute $60$}\\
u_{1} &= \res{820} &&\justif{$760+60$}
\end{calculs}
\explique{Chaque mois, il reste $95\,\%$ des trottinettes (on multiplie par $0{,}95$), puis
l'opérateur en ajoute $60$ : $u_{n+1}=0{,}95\,u_{n}+60$.}{$1-\dfrac{5}{100}=0{,}95$}
\cas{2.}
\begin{calculs}
v_{n+1} &= u_{n+1}-1\,200 &&\justif{définition de $v_{n}$, au rang $n+1$}\\
v_{n+1} &= 0{,}95\,u_{n}+60-1\,200 &&\justif{relation de récurrence}\\
v_{n+1} &= 0{,}95\,u_{n}-1\,140 &&\\
v_{n+1} &= 0{,}95(u_{n}-1\,200) &&\justif{$0{,}95\times 1\,200=1\,140$}\\
v_{n+1} &= 0{,}95\,v_{n} &&
\end{calculs}
\begin{calculs}
v_{0} &= u_{0}-1\,200 &&\\
v_{0} &= 800-1\,200 = \res{-400} &&
\end{calculs}
\explique{$(v_{n})$ est \res{géométrique, de raison $0{,}95$} et de premier terme
$v_{0}=-400$.}{forme $v_{n+1}=q\times v_{n}$}
\cas{3.}
\begin{calculs}
v_{n} &= -400\times 0{,}95^{n} &&\justif{$v_{0}\times q^{n}$}\\
u_{n} &= v_{n}+1\,200 &&\justif{car $v_{n}=u_{n}-1\,200$}\\
u_{n} &= 1\,200-400\times 0{,}95^{n} &&\justif{l'expression demandée}
\end{calculs}
\cas{4.}
\explique{$v_{0}<0$ et $0<q<1$ : le sens est inversé, $(v_{n})$ est strictement
croissante, donc $(u_{n})$ aussi : \res{strictement croissante}.}{propriété 7}
\explique{$0{,}95^{n}$ se rapproche de $0$ : on conjecture
$\res{\lim\limits_{n\to+\infty}u_{n}=1\,200}$.}{le parc se stabilise vers $1\,200$}
\cas{5.}
\explique{À la calculatrice (tableau de valeurs) : $u_{13}\approx 994{,}7$ et
$u_{14}\approx 1\,004{,}9$.}{premier rang où $u_{n}>1\,000$}
\explique{La ville disposera de plus de $1\,000$ trottinettes \res{au bout de $14$
mois}.}{la suite est croissante : elle reste ensuite au-dessus}
\end{exo}

\end{document}
